\chapter*{Who you gonna call?}
\addcontentsline{toc}{chapter}{Who you gonna call?}

Mein Name ist Anton. Meine Reise begann vor einiger Zeit in einem Auto, das man so auf den Straßen nicht alle Tage sieht: ein Mercedes W124, Baujahr 1993, ehemaliges Notarztfahrzeug. Mit seinem markanten Aufbau sieht er ein bisschen aus wie der Wagen der Ghostbusters. Und genau so fühlte es sich auch an: wie ein Abenteuer, bei dem man nie genau weiß, welches „Gespenst“ hinter der nächsten Kurve wartet. Der W124 war mein mobiles Zuhause, ein treuer Panzer aus Stahl, der mich durch verschiedenste Kulturen und Landschaften trug. Doch kein Land auf dieser Reise hat mich emotional so begeistert und gleichzeitig herausgefordert wie der Iran. Was ich hier erzähle, ist nur ein Bruchteil meiner Erlebnisse dort. Es ist die Geschichte vom spannendsten, aber gleichzeitig intensivsten Monat meines Lebens. Es war eine Zeit, an der die unendliche Hilfsbereitschaft der Menschen direkt auf die unnachgiebige Härte eines gnadenlosen Systems prallte. Ich weiß, die unten stehenden Zeilen sehen nach viel aus, aber bitte lest sie bis zum Schluss. Ich kann euch versprechen: Es wird sich lohnen. Ihr werdet verstehen, wie eine einfache Schachtel Kekse in einem gnadenlosen System plötzlich über Freiheit und Schicksal entscheiden kann.

\chapter*{Gehe nicht über Los}
\addcontentsline{toc}{chapter}{Gehe nicht über Los}

Mein Pass war zu diesem Zeitpunkt schon seit fast einer Woche weg. Ich saß im Ungewissen, aber ich blieb ruhig. Normalerweise bin ich ein optimistischer Mensch; ich löse Probleme, wenn sie auftauchen, anstatt mich von ihnen überwältigen zu lassen. Also gingen Payam und ich unvoreingenommen in das Gerichtsgebäude. „Das wird schon“, dachte ich mir. Es war bereits mein zweiter Besuch dort, aber dazu später mehr. Wir betraten das Zimmer des Gerichtsassistenten, der bis ins kleinste Detail geschniegelt war. Ein kurzer Wortwechsel, dann das Signal: Der Richter hat Zeit für uns. Wir wechselten in das Richterzimmer. Der Raum war trist: ein einfacher Holztisch, an dem der Richter thronte, und drumherum Metallstühle mit Ledersitzen, die zum Schutz vor Verschleiß sorgfältig in Plastikfolie eingeschlagen waren. Ich nahm an der Seite Platz, während Payam dem Richter direkt gegenüber saß. Die Stimmung im Raum war bedrückend. Payam hat eigentlich immer diesen freundlichen, entspannten Blick in den Augen, der einem sofort Vertrauen einflößt. Doch während die beiden sprachen, beobachtete ich, wie dieser Blick sich wandelte. Er wurde ernst, dann tief besorgt. Ich verstand kein einziges Wort des persischen Dialogs, aber das Unbehagen kroch mir langsam den Rücken hoch. Das Gespräch dauerte kaum zwei Minuten, doch sie fühlten sich an wie eine Ewigkeit. Als wir das Zimmer verließen, herrschte erst einmal Stille. Payam suchte nach Worten. Wahrscheinlich wurde ihm in diesem Moment klar, dass es für das, was er mir sagen musste, keine „richtige“ Formulierung gibt. „Anton“, sagte er schließlich schwerfällig, „sie werfen dir vor, dein Auto in den Iran geschmuggelt zu haben. Wir müssen in zwei Tagen wiederkommen. Mit einer Kaution von 48.000 \$. Wenn wir das Geld nicht haben, gehst du ins Gefängnis.“ Ich sah ihn an und fragte nur: „Ist das dein Ernst oder machst du Witze?“ Er sah mich an und bestätigte, dass das kein Witz war. 48.000\$. Eine Summe völlig jenseits meiner Realität. In meinem Kopf ploppte nur ein einziger, absurder Gedanke auf: Gehe nicht über Los. Ziehe keine 4000\$ ein. Geil! Wir gingen schweigend zum Auto und fuhren los.

\chapter*{Zwischen Raketen und rotem Sand}
\addcontentsline{toc}{chapter}{Zwischen Raketen und rotem Sand}

Nach drei Monaten Thailand-Aufenthalt kam ich Ende Februar 2025 zurück nach Tiflis, der Hauptstadt Georgiens. Ich holte Keti aus Armenien ab, die ich bereits aus Südostasien kannte, und wir verbrachten einige Zeit in Georgien, bis wir die nötigen Dokumente zusammen hatten. Die Visa für den Iran waren wirklich sehr schnell ausgestellt und stellten kein Problem dar. Die Ankunft des Carnet de Passage ließ allerdings auf sich warten. Dieses Dokument ist ähnlich zu einem Reisepass, bloß nicht für Menschen, sondern für Autos. Um ein solches zu bekommen, muss man abhängig vom Wert des Autos ein Deposit hinterlegen. In meinem Fall waren es 5.000 Euro. Als das Carnet endlich da war, entschieden wir uns recht zügig Richtung Iran zu fahren, da ich unbedingt die Nowruz-Feierlichkeiten miterleben wollte. So machten wir uns auf den Weg durch Armeniens atemberaubende Berglandschaften und kamen am 17. März im Iran an. An der Grenze wurde schnell klar, dass es nicht wirklich organisiert ist. Das Carnet wurde bereits gestempelt, aber ich sollte noch zu einem Herrn in einen Container. Dieser Herr telefonierte rum, schien aber scheinbar nicht die nötige Information zu bekommen. Er wirkte sehr entspannt und gleichgültig. Nach kurzer Zeit ließ er mich einfach gehen. Mittlerweile wurde mir bewusst, was sein Problem war: Laut Gesetz ist er eigentlich dazu verpflichtet, mir ein Ausreisedatum für das Auto zu geben. Ihm war scheinbar nicht bewusst, wann das ist. Mir persönlich natürlich auch nicht. Allgemein ist es nur sehr schwer, als Ausländer Informationen aus dem Iran zu beziehen. Das liegt nicht nur an der Sprachbarriere, sondern auch daran, dass sämtliche iranische Webseiten nur aus dem Iran verfügbar sind. In diesem Moment habe ich nicht verstanden, was genau das Problem dieses Mannes war. Es war mir auch egal; ich war froh, dass er uns einfach fahren ließ. Der Iran war aufregend und anders als alles andere, wo ich bisher war. Wir fuhren in den nächsten Wochen über Tabriz, Rascht, Teheran und Isfahan nach Shiraz. Nach circa einem Monat gemeinsamer Reise musste Keti wieder abreisen. Ich führte meine Tour allein fort und verbrachte einen weiteren Monat im Süden des Landes. Da der Liter Benzin gerade mal 2 Cent kostet, nutzte ich das natürlich aus und besuchte zahlreiche Orte und Inseln. Meine absolute Empfehlung ist die Insel Hormuz. Ich habe so eine Art Natur noch nie gesehen: Vom Boden aus gelbem Kristallgestein über tiefroten bis hin zu funkelndem schwarzen Sand hat diese Insel wirklich viel zu bieten. Von Hormuz ging es für mich zurück auf die Insel Qeshm, wo ich Alen kennenlernte, einen Reisenden aus Kroatien, der mit seinem VW T4 bereits seit über neun Jahren reiste. Er war auf der Insel gestrandet, weil er ein größeres Problem mit seinem Wagen hatte. Da ich zuvor zwei Monate lang gar keinen Alkohol getrunken hatte, kam mir die Abwechslung gelegen: Alen hatte eine Connection zu einem Arak-Händler. Wir tranken den selbstgebrannten Schnaps zusammen und redeten über unsere Pläne. Alen erwähnte, dass er bereits vorher zwei Mal seinen Van über einen längeren Zeitraum im Iran gelassen hatte. Es sei kein Problem, solange das Carnet noch gültig ist. Rückblickend betrachtet ist die Lektion klar: Traue keinem betrunkenen Kroaten, wenn es um komplexe gesetzliche Bestimmungen geht. Zu diesem Zeitpunkt machte ich mir jedoch keinerlei Gedanken dazu, da ich auf dem Weg nach Nepal war. Am 7. Mai war es für uns soweit, die Insel zu verlassen. Mit einem Abschleppseil zog ich Alens T4 Richtung Fähre. Wir fuhren in der sengenden Hitze bei 40 Grad mit 30 km/h über die Straßen. Mein W124 hat keine Klimaanlage, woran ich mich bei diesen Temperaturen zwar gewöhnt hatte, was die Fahrt aber nicht gerade entspannter machte. Auf der Fähre saßen wir mit den Mitarbeitern zusammen und tranken Tee. Ich erzählte, dass ich auf dem Weg über Pakistan und Indien nach Nepal sei. Einer der Männer machte mit seinen Händen Raketenbewegungen von Indien nach Pakistan. Ich dachte, er meinte die Scharmützel im Kaschmir-Konflikt, die mal wieder am Eskalieren waren und ging davon aus, dass diese sich wie auch vorher schon beruhigen würden. Ich lag falsch. Als wir am Festland ankamen, schaute ich in die Nachrichten. Dort stand es schwarz auf weiß: Indien hatte an diesem Tag Pakistan mit Raketen angegriffen. Das war eine deutlich größere Sache als die paar kleinen Scharmützel vorher. Mir wurde bewusst, dass die Grenze zwischen Indien und Pakistan geschlossen bleiben wird. Mein Plan ging nicht auf. Ich beschloss, einen kühlen Kopf zu behalten, und mietete mir ein Hotelzimmer in Bandar Abbas, um online zu recherchieren. Dort fand ich in zahlreichen Reiseforen und Blogs Bestätigungen für Alens Aussage. Es gab nun zwei Möglichkeiten: entweder ich fahre über Pakistan nach Afghanistan und weiter Richtung Tadschikistan oder ich lasse den Wagen im Iran und fliege zurück nach Südostasien. Letzteres wirkte deutlich attraktiver, als mich in Afghanistan mit den Taliban rumzuärgern. Ich beschloss, eine Nacht darüber zu schlafen. Am nächsten Morgen buchte ich mein Ticket für den 12. Mai von Shiraz nach Bangkok und machte mich auf den Weg nach Shiraz. Dort angekommen lernte ich Razool kennen, Alens Mechaniker. Razool ist ein wirklich sehr netter Mann, dessen Herz für VWs schlägt. Er bot mir an, meinen Wagen bei ihm auf seinem privaten Grundstück zu parken. Ich sah das als die beste Option an. An meinem Abreisetag am 12. Mai gab es jedoch noch ein unerwartetes Hindernis. In der großen Hitze der letzten Wochen war es mir egal geworden, dass man wegen der Scharia eigentlich keine kurzen Hosen tragen sollte. Für Ausländer gelten im Iran allgmein etwas andere Gesetze. Da ich nach Bangkok flog, hatte ich keine einzige lange Hose dabei. Am Flughafen sagte mir ein Beamter, ich könne so nicht fliegen. Er durchsuchte mein Gepäck, fand aber natürlich nichts. Ich schlug ihm vor, mein großes Decathlon-Handtuch wie einen Rock über die Beine zu binden. Er stimmte zu. So saß ich mit meinem umgebundenen Handtuch in der Wartehalle und bestieg schließlich das Flugzeug.

\chapter*{Zwischen den Welten}
\addcontentsline{toc}{chapter}{Zwischen den Welten}

Die folgenden Monate verbrachte ich in Thailand, Laos, Vietnam, China und Japan. Ich war weit weg von den bürokratischen Hürden des Nahen Ostens. Während ich diese Länder bereiste, stand mein Mercedes weiterhin sicher geglaubt in Shiraz auf einem privaten Grundstück. Doch die geopolitische Lage im Iran änderte sich im Juni 2025 schlagartig durch den 12 Tage Krieg. Schwere Luftangriffe durch Israel und die USA trafen das Land. Mitten in dieser Zeit der Instabilität meldete sich Razool bei mir. Er bat um Fotos meines Reisepasses und des Carnet de Passage. Er begründete dies damit, dass die Polizei nachfragen könnte. Wie ich später erfuhr, hatte Razool sich beim iranischen Geheimdienst gemeldet, weil mein ausländisches Fahrzeug auf seinem Grundstück geparkt war. In einer Zeit höchster militärischer Paranoia war der Wagen für die Behörden so zumindest offiziell bekannt und nicht versteckt. Parallel dazu änderte sich auch die Situation für Alen. Er hatte seinen T4 ebenfalls im Iran zurücklassen müssen, da eine Reparatur vor Ort nicht möglich war. Er informierte mich darüber, dass das Zurücklassen des Fahrzeugs entgegen seiner ersten Vermutung doch zu massiven Problemen führen könnte. Alen selbst hatte bei seiner Fristüberschreitung, die parallel zu meiner verlief, sehr viel Glück gehabt und lediglich etwa 70 Euro bezahlt. Ich begann daraufhin selbst intensiver zu recherchieren. Ich stieß auf Berichte, nach denen eine Geldstrafe von etwa 1000 Euro für die Zeitüberschreitung fällig würde. Für mich war das ein kalkulierbares Risiko und eher eine hohe Gebühr für meine gewonnene Reisefreiheit in Südostasien. Da ich wegen der Fristüberschreitung ohnehin schon in der rechtlichen Klemme steckte, sah ich keinen Grund für eine überstürzte Rückkehr. Ich entschied mich zu warten, bis sich die militärische Lage im Land etwas beruhigt hatte. Kurz vor meiner geplanten Rückreise gab es noch einmal Kontakt zu Razool. Er half mir dabei, die Lichtmaschine an meinem Mercedes zu reparieren, damit der Wagen für die kommende Fahrt einsatzbereit war. Am 30. Dezember 2025 landete ich schließlich am Imam Khomeini Flughafen in Teheran. Die Einreise verlief ohne große Probleme. Dennoch war die Atmosphäre bei meiner Ankunft angespannt. Mir war zu diesem Zeitpunkt noch nicht bewusst, wie sehr sich die Lage im Land zugespitzt hatte. Ich kehrte in ein Land zurück, das nach dem Krieg im Sommer nun im Inneren brannte.

\chapter*{Im Vakuum des Blackouts}
\addcontentsline{toc}{chapter}{Im Vakuum des Blackouts}

Nach meiner Landung verbrachte ich zunächst ein paar Tage in Teheran, bevor ich mit dem Nachtzug nach Shiraz fuhr. Ich freute mich sehr auf die Rückkehr an diesen Ort, der für mich zu den liebsten im ganzen Land gehört. Das lag zum einen an der Freundschaft zum Gastgeber meines Hotels, zum anderen daran, dass ich dort endlich meinen Mercedes wieder in Empfang nehmen konnte. Jalal, seine Frau Zara und ihr Bruder Ali empfingen mich sehr herzlich in ihrem Zuhause. Es ist ein wunderschönes Haus im alten iranischen Stil mit bunten Fenstern, die zum grünen Innenhof ragen. Aufgrund der landesweiten Proteste blieb ich die meiste Zeit im Haus und ging fast nur zum Essen vor die Tür. Es herrschte eine familiäre Stimmung und ich verbrachte viel Zeit mit Gesprächen mit den anderen Gästen. Da Razool für einige Tage auf einem VW Treffen in Buschehr war, musste ich etwas auf meinen Wagen warten. Als ich ihn schließlich abholte, war er unter einer dicken Staubschicht begraben. Ich setzte mich hinein und drehte den Schlüssel um. Der Wagen sprang sofort an, genau wie ich es gewohnt war. Auch wenn der Mercedes von außen und innen etwas heruntergerockt aussieht, läuft der Motor wie eine eins. Der Rost an einigen Stellen war über die Monate allerdings deutlich sichtbarer geworden. Da Reparaturen im Iran im Vergleich zu Deutschland extrem günstig sind, suchte ich über Razool nach einem Lackierer. Wir einigten uns auf einen Preis von etwa 200 Euro. Für diese Summe lackierte er die Heckklappe, die Motorhaube und mehrere kleine Stellen rund um das Auto. Außerdem flexte er eine alte Außensteckdose ab und schweißte ein neues Blech ein. In Deutschland hätte ich für diese viertägige Arbeit weit über 1000 Euro bezahlt. Die Wartezeit für die Lackierung war kein Problem, da ich ohnehin noch einmal nach Teheran musste. Mein russischer Pass lag beim Konsulat, das wegen des Neujahrsfestes bis zum 12. Januar geschlossen blieb. Bevor ich den Wagen beim Lackierer abgab, suchte ich stundenlang nach Winterreifen. Meine geplante Route sollte mich durch die Berge des kurdischen Teils Irans und später nach Armenien führen. Ich machte die Erfahrung, dass Sommerreifen dort definitiv keinen Spaß machen. Doch die Suche blieb erfolglos. Die Leute verstanden oft gar nicht, was ich suchte, oder boten mir Schneeketten an, die jedoch ebenfalls nicht verfügbar waren. Ich gab schließlich auf und brachte den Mercedes zum Lackierer. Am Mittwoch, den 7. Januar, zeigte sich die Anspannung im Land deutlich. Ich wollte etwas zu essen finden, doch obwohl es ein normaler Wochentag war, hatten alle Geschäfte geschlossen. Es war ein beängstigendes Bild, aber gleichzeitig spannend zu sehen, wie stark die Bevölkerung durch diesen Boykott gegen das Regime zusammenhielt. Der wahre Schrecken folgte am nächsten Tag, als das Internet im gesamten Land einfach ausgeschaltet wurde. Zuerst dachte ich an eine kurze Unterbrechung über das Wochenende, doch ich irrte mich. Da nur noch nationale Webseiten funktionierten, konnte Jalal immerhin Zugtickets nach Teheran für mich buchen. Nachdem ich meinen Wagen am Freitag in tadellos lackiertem Zustand abgeholt hatte, machte ich mich am Sonntagabend mit dem Nachtzug auf den Weg in die Hauptstadt. Am Montagmorgen kam ich in Teheran an und war der erste in der Schlange vor dem Konsulat. Dort teilte man mir mit, dass mein Pass zwar vorlag, ich ihn aber wegen des Blackouts eigentlich nicht erhalten könne, da die Abholung nach Moskau gemeldet werden müsse. Nach einer langen Diskussion hatte ich Glück, da ich zusätzlich mit meinem deutschen Pass reiste. Sie händigten mir das Dokument schließlich aus, warnten mich aber, dass der Pass im System erst aktiviert würde, sobald wieder Internet bestünde. Ich war erleichtert, denn ein neuer Pass wäre ein riesiger Aufwand gewesen und ich musste den Iran schnellstmöglich verlassen. Vor der Rückreise suchte ich noch die deutsche Botschaft auf. Ich hoffte verzweifelt auf einen Weg, meine Eltern zu erreichen, die zwar von meinem Aufenthalt im Iran wussten, aber vom plötzlichen Blackout völlig unvorhergesehen getroffen wurden. Es gab jedoch keine Möglichkeit für eine Verbindung nach außen. Zurück in Shiraz verbrachte ich eine letzte Nacht bei Jalal als sein einziger Gast. Am nächsten Tag begann die tagelange Fahrt durch die malerische Berglandschaft des iranischen Westens Richtung Täbris. Da ich zumindest eine iranische Navigationsapp hatte, kam ich gut voran. Das erzwungene digitale Detox fühlte sich anfangs sogar nach Freiheit an, da ich endlich wieder Stunden mit meiner Lieblingsbeschäftigung, dem Autofahren, verbringen konnte. Dennoch machten mir die Kontrollen zu schaffen: Ich wurde an fast jeder Station angehalten, mein Auto wurde durchsucht und einmal hielt man mich sogar für ein zweistündiges Interview fest. Mit jedem Kilometer Richtung Norden wuchs die Sorge um meine Liebsten daheim, während die Grenze zu Armenien langsam näher rückte.

\chapter*{Eisiger Stillstand}
\addcontentsline{toc}{chapter}{Eisiger Stillstand}

Am Dienstag, den 20. Januar 2026, wachte ich in meinem Hotel in Täbris auf, frühstückte in Ruhe und machte mich schließlich bereit für die Fahrt nach Norduz an die Grenze. Schon kurz nach dem Aufbruch deutete sich an, dass dieser Tag nicht unbedingt mein bester werden würde. Die Temperaturen waren im Vergleich zu den vergangenen Tagen spürbar gesunken und dichter Schneefall setzte ein. Obwohl ich im Vorfeld extra Scheibenwischflüssigkeit mit Frostschutz gekauft hatte, erwies sich das Mittel als völlig wirkungslos; es schien kaum mehr als blau gefärbtes Wasser zu sein. Durch den aufgewirbelten Schneematsch und den Dreck der Straße war meine Windschutzscheibe bereits nach wenigen Minuten so stark verschmutzt, dass ich kaum noch etwas erkennen konnte. Zum Glück verfügte mein W124 über ein beheiztes Wischwassersystem. Anfangs half die Wärme der Flüssigkeit sogar dabei, den durch die Kälte bereits leicht verbogenen Scheibenwischer wieder in seine ursprüngliche, gebogene Form zu bringen. Doch getreu dem Motto „viel hilft viel“ setzte ich immer größere Mengen des beheizten Wassers ein. Das erwies sich als fataler Fehler: Die Flüssigkeit blieb am Wischerblatt hängen, und unter dem Einfluss des eisigen Fahrtwinds fror das Ganze auf der Autobahn schlagartig ein. Der zuvor elastische und gebogene Wischer verwandelte sich in einen starren, geraden Stab, der keinen Kontakt mehr zur Scheibe hatte. Meine Sicht war dadurch fast vollständig blockiert. Mit Müh und Not rettete ich mich zu einer Tankstelle, an die ein kleiner Kiosk angeschlossen war. Ich fragte den Besitzer nach einem vernünftigen Frostschutzmittel und er hatte glücklicherweise ein passendes Produkt vorrätig. Ich kippte die Flüssigkeit direkt über den Scheibenwischer, doch die zentimeterdicke Eisschicht zeigte sich davon völlig unbeeindruckt. Als ich den Kioskbesitzer um Rat fragte, hatte dieser eine pragmatische Lösung parat: Er füllte kochendes Wasser aus seinem Samowar in einen Kessel und half mir dabei, die Mechanik eigenhändig aufzutauen. Ich war erleichtert und froh, meine Fahrt endlich mit klarer Sicht fortsetzen zu können. Wenig später folgte jedoch die nächste Überraschung. Ich war seit Schiras rund 1.600 Kilometer bis nach Täbris gefahren und hatte dabei fast durchgehend verschneite Landschaften durchquert, in denen die Straßen jedoch immer tadellos geräumt waren. Warum ich ausgerechnet jetzt auf einen völlig zugeschneiten Autobahnabschnitt geriet, war mir ein absolutes Rätsel. Mit meinen Sommerreifen schlich ich mit gerade einmal 20 km/h über den Asphalt. Es dauerte etwa eine Stunde, bis sich die Straßenverhältnisse wieder besserten. Gezeichnet von den Strapazen dieser Fahrt erreichte ich schließlich nach insgesamt vier Stunden die Grenze in Norduz. Bei meiner Ankunft hatte ich das Gefühl, einigermaßen vorbereitet zu sein. Ich hatte einen Brief aufgesetzt, in dem ich meine Lage beschrieb und darum bat, eine mögliche Geldstrafe gering zu halten. Mein Plan war es, ein Dokument zu erhalten, das bescheinigte, dass mit meinem Auto alles in Ordnung sei. Mit etwas Glück wollte ich danach sogar in Richtung der türkischen Grenze weiterfahren, da dies mein eigentliches Ziel war, um den Rückweg nach Europa anzutreten. Ich parkte den Mercedes vor der Grenze, stieg aus und orientierte mich kurz. In diesem Moment kam ein Mann auf mich zu, etwa 1,80 Meter groß, korpulent, mit einem freundlichen Lächeln. Er wirkte auf mich wie ein Händler und sprach mich in einem für iranische Verhältnisse ungewöhnlich guten Englisch an. Ich erklärte ihm meine Situation und dass ich dieses Dokument benötigte. Er wies mir den Weg, woraufhin ich erst einmal zu Fuß die erste Schranke passierte. Auf einem kleinen Zettel hatte ich mir die richtige Abteilung notieren lassen, die ich nun suchte. Schließlich fand ich das richtige Büro und übergab einem Mitarbeiter sowohl meinen Notizzettel als auch den vorbereiteten Brief. Er verschwand damit in einem Büro. Nach einer Weile erschien ein anderer Beamter, der offenbar zuständig war. Er war klein, hatte vorstehende Zähne und starrte mich durch seine dicke Brille mit schielenden Glubschaugen an; sein Pullover war rundherum mit Schuppen bedeckt. Der Einfachheit halber nennen wir diesen Mann in dieser Geschichte "Brille". Er lächelte mich nett an und sprach ein gebrochenes Englisch, wobei es beim Reden so wirkte, als würde er auf seiner eigenen Zunge herumkauen. Er erklärte mir, dass ich den Wagen erst durch die erste Schranke bringen müsse, damit sie mir das Dokument ausstellen könnten. Da ich keine andere Möglichkeit sah, holte ich das Auto nach. Brille ging nun mit mir in ein paar andere Büros. Die Beamten dort überflogen meinen Brief eher gleichgültig und redeten dabei auf Farsi. Ein Begriff kristallisierte sich jedoch heraus, der mir auffiel, auch wenn ich ihn nicht verstand: Tasirat. Im Büro des Zollchefs wurden schließlich mein Carnet de Passage und mein Pass einkassiert. Brille händigte mir eine Kopie aus, auf der er vermerkte, dass er meinen Pass habe, und seine Telefonnummer notierte. Am Ende prüften wir noch, ob die Karosserienummer des Wagens mit der Nummer auf dem Carnet übereinstimmte. Ich fragte ihn im Anschluss, ob es ein Problem geben würde, und erwähnte, dass ich das Gefühl hätte, der Erste mit solch einem Problem zu sein, da wir zuvor mit so vielen Leuten geredet hatten. Er antwortete einfach nur mit „No, no problem!“ und lächelte. Ich dachte mir: Okay, dann ist ja alles gut. Er sagte mir, ich solle ins Hotel gehen und am nächsten Morgen wiederkommen, um zum Gericht zu gehen. Ich ging davon aus, dass das Gericht lediglich die Strafe für die Fristüberschreitung berechnen würde und ich diese dann dort bezahlen könnte. Ich war nicht sehr glücklich über die Situation, sah es jetzt aber auch nicht als großes existenzielles Problem an. Mit meiner Passkopie versuchte ich, im Hotel vor Ort einzuchecken. Der Mitarbeiter wusste jedoch nicht, ob er mich aufnehmen dürfe, da ich keinen Originalpass hatte und das zuständige Büro gerade zu war. Er bat mich, in einer Stunde wiederzukommen. In diesem Zusammenhang muss man erwähnen, dass man beim Einchecken im Iran seinen Pass im Hotel abgeben muss, solange man dort bleibt. Von den 42 Ländern, die ich bereist habe, ist es das einzige mit dieser seltsamen Regel, die sowohl für Ausländer als auch für Iraner gilt. Ich ging erst einmal ins Restaurant nebenan, um die Wartezeit mit einer Suppe und Kebab zu überbrücken. Wenig später betrat auch der Mann, den ich kurz zuvor an der Grenze getroffen hatte, den Raum und setzte sich zu mir. Während wir aßen, erklärte ich ihm meine Lage, woraufhin er sich ein wenig über meine Gutgläubigkeit lustig machte; er meinte, die Zöllner hätten mich wohl mit einem Vorwand dazu gebracht, mein Auto freiwillig auf das abgesperrte Gelände zu fahren. Ich dachte mir dabei nicht viel. Er erzählte mir, dass er einige Jahre in Australien und Dubai gelebt hatte, was sein fließendes Englisch erklärte. Dass er ausgerechnet hier an der Grenze Zeit verbrachte, hatte einen pragmatischen Grund: Er besaß eine armenische SIM-Karte, die sich hier im Grenzgebiet ins Nachbarnetz einwählte. Inmitten des landesweiten Blackouts war dies der einzige Ort weit und breit mit funktionierendem Internet. Er war so hilfsbereit, mir sein Handy zu überlassen, damit ich meinen Eltern endlich ein Lebenszeichen geben konnte. Er stellte sich als Payam vor und ich speicherte seine Nummer direkt in meinem Handy. Nach dem Essen begleitete er mich zurück zum Hotel und half mir beim Einchecken, wo die Angestellten dank seiner Vermittlung die Passkopie schlussendlich doch akzeptierten. Ich verabschiedete mich von ihm in dem festen Glauben, meine Reise am nächsten Tag ohne große Umschweife fortsetzen zu können.

\chapter*{Das versiegelte Schicksal}
\addcontentsline{toc}{chapter}{Das versiegelte Schicksal}

Am nächsten Morgen ging ich wieder an die Grenze, um mich mit Brille zu treffen. Er erklärte mir direkt, dass der Richter heute nicht im Haus sei und ich bis Samstag warten müsse. Um neun Uhr morgens sollte ich wieder erscheinen, damit er gemeinsam mit mir zum Gericht fahren könne. Er machte einen freundlichen Eindruck und sprach ein wenig Englisch, was mich in dem Glauben ließ, dass alles seinen geregelten Gang gehen würde. Es war zwar keine ideale Situation, aber ich hatte keine andere Wahl, als zuzustimmen. Beim Verlassen des Grenzgeländes traf ich zufällig Payam und erzählte ihm von dem Aufschub. Er fragte mich, ob wir zusammen noch einmal hineingehen sollten, um mit Brille zu reden. Da ich zu diesem Zeitpunkt noch immer davon ausging, dass ich lediglich zum Gericht müsse, um die Strafe für die Zeitüberschreitung des Autos zu bezahlen, lehnte ich ab. Ich wollte keine Umstände machen. Das war ein großer Fehler, wie sich später noch herausstellen sollte. Wir gingen mit Payam Sandwich essen und unterhielten uns über Gott und die Welt. Ich speicherte über WhatsApp auch Payams armenische Nummer. Danach suchten wir nach einer armenischen SIM-Karte für mich. Es wäre einfach gewesen, eine solche zu bekommen, wenn ich das Land hätte verlassen können, aber ohne Pass und mit dem One-Way-Visum war das nicht möglich. Nach einigem Suchen fanden wir einen Mann, der voraktivierte SIM-Karten verkaufte. Die zweite Karte, die wir testeten, funktionierte gut. Ich freute mich, nach knappen zwei Wochen endlich wieder Internet zu haben. Wir besorgten noch einen Gas-Campingkocher für mich, damit ich mir im Hotel etwas zubereiten konnte. Danach verabschiedeten wir uns und ich ging zurück ins Hotel. Die nächsten zwei Tage verliefen eher unspektakulär. Ich nutzte die Zeit, um die Nachrichten nachzuholen, die ich während des Blackouts verpasst hatte. Dabei erfuhr ich von dem blutigen Massaker, das die Regierung im Schatten der digitalen Dunkelheit an Demonstranten verübt hatte. Nebenbei arbeitete ich an der Webseite für mein zukünftiges Unternehmen. Am Freitag wollte ich Payam anrufen, um mich für den nächsten Tag abzustimmen, doch seine Nummer war plötzlich unauffindbar. Ich konnte mir das absolut nicht erklären, da ich sie beim Speichern über WhatsApp noch gesehen hatte. Am Samstagmorgen stand ich um neun Uhr pünktlich an der Grenze, aber von Brille war keine Spur. Nach einer Stunde fragte ich eine Frau in seiner Abteilung nach ihm. Sie rief ihn an, aber da sie kein Englisch sprach, gab es keine wirkliche Kommunikation. Sie druckte einen Brief aus, versiegelte ihn, schrieb eine Adresse auf den Umschlag, zeigte mit Ihrem Finger darauf darauf und sagte nur: "Taxi". Ich suchte mir draußen einen Fahrer, der mich für umgerechnet 7 Euro zum etwa 50 Minuten entfernten Gericht in Jolfa brachte. Der Preis war für iranische Verhältnisse teuer, aber ich hatte keine andere Wahl. Am Gericht angekommen, schnappte sich der Taxifahrer den Brief und ging mit mir hinein. Ich hatte eigentlich damit gerechnet, dass Brille dort schon auf mich warten würde, aber er war nicht da. Der Taxifahrer ging zuerst zum Assistenten des Richters. Wir mussten warten. Aufgrund der Aufregung musste ich zur Toilette und tippte in meinen Handyübersetzer, ob es hier eine gebe. Der Assistent antwortete mir nicht. Stattdessen machte er sich scheinbar über mich lustig, indem er etwas zu den anderen Besuchern sagte und lachte. Ein wirklich toller Ort. Nach zehn Minuten hatte der Richter Zeit für uns. Der Taxifahrer und ich gingen hinein und die beiden unterhielten sich auf Farsi, während ich nur danebenstehen konnte. Der Taxifahrer machte eine Handbewegung, dass wir los konnten, und zeigte mir beim Rausgehen mit seinen Fingern eine Drei. Ich verstand nicht, ob er drei Stunden oder drei Tage meinte. Draußen gab ich ihm mein Handy, damit er mit dem Übersetzer schreiben konnte, was los sei. In diesem Moment traf mich die Realität wie ein Schlag: Der Taxifahrer war einfach nicht in der Lage zu schreiben. Ich starrte ihn an und verstand gar nicht, was gerade passiert war. Da stand ich nun, wusste absolut nicht, was juristisch auf mich zukam, und meine einzige Brücke zu diesem Justizapparat war ein Mann, der keine einzige Taste auf meinem Handy bedienen konnte. Es war ein absurder, fast schon beängstigender Moment der totalen Isolation. Draußen standen ein paar andere Männer. Ich fragte diese mit Handgesten, ob sie aufschreiben könnten, was der Taxifahrer sagte. Einer tippte in mein Handy und ich las in der Übersetzungsapp: "3 Tage". Das war eine Katastrophe, da in drei Tagen der letzte Tag meines Visums wäre. Ich ging zurück zum Richter und zeigte ihm über den Übersetzer, dass mein Visum in drei Tagen abläuft. Ich machte Gesten, dass er sich über den Übersetzer mit mir verständigen könne, doch er machte nur eine Handbewegung, dass er mein Handy nicht nehmen wollte. Da mein Handy im Iran gesperrt war, konnte ich selbst niemanden anrufen. Ich zeigte auf sein Handy und machte eine Telefongeste, doch er lehnte auch das ab. Ich ging raus und fand einen Mann, der mir erlaubte, mit seinem Smartphone zu telefonieren. Ich rief Jalal aus Shiraz an, den ich bereits vorgewarnt hatte. Ich wollte dem Richter das Telefon geben, doch er lehnte wieder ab. Der Assistent war gerade mit im Raum und ging zurück an seinen Arbeitsplatz. Ich folgte ihm und gab ihm das Handy. Er nahm es, war aber sichtlich nicht gewillt, mit Jalal zu reden. Er gab mir das Telefon zurück und sagte zu den anderen Anwesenden mit einem Kopfschütteln etwas von "Hotel Shiraz". Jalal sagte mir am Telefon, dass der Mann super unhöflich war. Er erzählte mir, dass sie behaupteten, mir helfen zu wollen, damit ich nicht ins Gefängnis müsse. Jalal nahm das in dem Moment jedoch so wahr, als wäre das nur Geschwafel und die Situation gar nicht so schlimm. Der Assistent drückte dem Taxifahrer nun einen Brief in die Hand, und dieser unterschrieb den Erhalt dessen. Wir machten uns auf den Rückweg nach Norduz. Dort angekommen, winkte der Taxifahrer dem Grenzbeamten mit dem Brief, als sei er extrem wichtig. Ich fühlte nur Verachtung für diesen Mann, der scheinbar keinerlei Gefühl dafür hatte, in was für eine Lage er mich gebracht hatte. Wir gingen ins Gebäude und er gab einem Mitarbeiter den Brief. Es ging etwas hin und her, bis ein Verantwortlicher bestimmt wurde. Der Mann teilte mir mit, ich solle in einer Stunde wiederkommen. Kaum hatte ich die Grenze verlassen, bekam ich einen Anruf. Ich schaue aufs Handy und sehe Payams armenischen WhatsApp-Namen. Ich weiß nicht, ob das nur mir so geht, aber wenn ich bei WhatsApp eine Nummer zu einem bestehenden Kontakt hinzufüge, wird der Name oft einfach durch den Profilnamen des Kontakts überschrieben. Das war mir bereits früher aufgefallen, in dieser stressigen Zeit habe ich das aber schlicht vergessen. Ich freute mich und erzählte ihm vom Briefwechsel. Nach dem Essen ging ich zurück zum Zoll. Der Beamte schrieb einen Brief fertig und brachte ihn mit mir zusammen zu seinem Vorgesetzten. Auf dessen Tisch sah ich in dem Dokument eine sehr lange Zahl. Ich holte mein Handy raus, rechnete die Zahl um und dachte, das seien 1.600\$ Strafe. Der Vorgesetzte rechnete jedoch noch einmal nach und drehte mir den Taschenrechner hin. Ich hatte eine Null übersehen. 16.000\$. Mir blieb das Herz stehen und mir klappte buchstäblich die Kinnlade runter. Der Zoll-Chef nickte die Zahl auf Anfrage einfach ab, als wäre es das Normalste der Welt. Ich war wie gelähmt vor Schock. Der Vorgesetzte faltete den Brief, versiegelte den Umschlag und sagte nur ein einziges Wort: „Problem“. Er gab mir zu verstehen, dass ich diesen versiegelten Brief am nächsten Tag wieder zum Gericht nach Jolfa bringen müsse. Ich musste mich erst einmal kurz fangen, bevor ich in mein Handy tippte, dass ich eine Kopie dieses Briefes brauche, weil es schließlich um mich ging und ich jetzt wohl einen Anwalt benötigen würde. Er überlegte einen Moment, druckte den Brief dann aber tatsächlich aus. Er riss jedoch fast alles davon ab und ließ nur den schmalen Streifen mit der Zahl stehen. Eine andere Nummer, vermutlich die Prozessnummer, machte er zusätzlich mit einem Kugelschreiber unkenntlich. Ich nahm diesen Papierfetzen und verließ die Grenze. Ich wusste nicht mehr, wo oben und unten war. Es war ein Schock, den ich in dieser Form bisher nicht kannte, da ich normalerweise selbst in stressigen Situationen sehr fokussiert bleibe. Hier jedoch fühlte ich mich einfach nur vollkommen machtlos. In diesem Zustand rief ich Payam an und schilderte ihm alles. Er bot mir sofort an, zu ihm nach Täbris zu kommen, dort zu übernachten und am nächsten Tag gemeinsam zum Gericht zu fahren. Ich nahm die Einladung an und machte mich mit einem Taxi auf den Weg nach Täbris.

\chapter*{Der 16.000-Dollar-Panzer}
\addcontentsline{toc}{chapter}{Der 16.000-Dollar-Panzer}

Die Fahrt nach Täbris dauerte etwa drei Stunden. Auf halbem Weg wechselte ich vom Taxi in ein Sammeltaxi, in dem ein Iraner aus Dubai saß, der fließend Englisch sprach. Ich erzählte ihm und dem Taxifahrer meine gesamte Geschichte und zeigte ihm den Schnipsel, auf dem die vermeintliche Strafe stand. Der Mann klärte mich jedoch schnell auf: Die 16.000 Dollar waren nicht meine Strafe, sondern der vom Zoll festgesetzte Wert meines Autos. Diese Information beruhigte mich ein wenig. Ich zeigte ihm ein Foto meines Wagens, und wir scherzten gemeinsam darüber, dass dieser alte Mercedes tatsächlich mit einer solchen Summe bewertet wurde. In Täbris angekommen, wurde ich von Payam und seinem Freund Saleh empfangen. Saleh arbeitete als Manager am Flughafen von Täbris und verfügte daher über sehr gute Kontakte in der Stadt und der gesamten Provinz Ost-Aserbaidschan, da Täbris die Hauptstadt dieser Provinz ist, in der auch Norduz liegt. Wir fuhren in eine Mall, wo Saleh mehrere Telefonate führte, um eine Lösung zu finden. Da sich vorerst nichts Vielversprechendes ergab, fuhren wir zu Payam nach Hause. Er war so gastfreundlich, mir sein eigenes Zimmer zu überlassen, während er selbst auf der Couch im Wohnzimmer schlief. Am nächsten Morgen machten wir uns auf den Weg zum Gericht in Jolfa. Es war gar nicht so leicht, das richtige Gebäude zu finden. Unterwegs fragten wir einen Passanten nach dem Weg, der den Ort anscheinend sehr gut kannte. Er erzählte uns, dass er bereits 15 Fälle bei meinem zuständigen Richter gehabt hatte und dieser ein ziemliches Arschloch sei. Diese Einschätzung deckte sich mit dem, was ich zuvor schon mehrfach von anderen Leuten an der Grenze gehört hatte. Ich fragte mich in diesem Moment ernsthaft, warum ich ausgerechnet an diesen Mann geraten musste. Am richtigen Gericht angekommen, teilte uns der Assistent jedoch mit, dass der Richter nicht anwesend war. Es war also ein weiterer verlorener Tag, an den ich mich langsam schon gewöhnte. Wir beschlossen, zurück nach Norduz zu fahren, um den überzogenen Wert des Autos offiziell in Frage zu stellen. Zuerst sah es gar nicht schlecht aus: Wir fragten uns von Büro zu Büro durch, bis der Chef des zuständigen Mitarbeiters schließlich eine handschriftliche Notiz auf einen Zettel schrieb. Mit diesem Zettel gingen wir zurück zu dem Mitarbeiter, der den Wert festgesetzt hatte. Dieser erklärte uns dann aber unmissverständlich, dass der Wert fest sei und er diesen Wert bereits mit Wohlwollen möglichst klein gewählt hatte. Wenn wir diesen anfechten wollten, müssten wir das über Teheran regeln, was jedoch zwei bis drei Wochen dauern würde. Im Nachhinein erklärte mir Payam auch die absurde Logik hinter der Wertermittlung: Da der Beamte mein Auto nicht im System finden konnte, nahm er einfach einen Mercedes aus dem Jahr 2018 als Grundlage. Er rechnete 25 Jahre Wertverlust ab und erhob zusätzlich eine Luxussteuer von 70 \%. Es war ein zutiefst professionelles Vorgehen, genau wie man es sich bei einem Grenzbeamten in Norduz vorstellt. Da wir nichts weiter ausrichten konnten, verbrachten wir die Nacht in dem mir bereits bekannten Hotel am Grenzgebiet.

\chapter*{No Problem, it's just 18.000€}
\addcontentsline{toc}{chapter}{No Problem, it's just 18.000€}

Am nächsten Morgen machten wir uns erneut auf den Weg nach Jolfa ins Gericht. Der geschniegelte Assistent machte den Eindruck, als ob er bereits auf uns gewartet hätte. „Der Richter hat jetzt Zeit für euch“, signalisierte er knapp. Wir betraten das Richterzimmer. Ich nahm an der Seite Platz, während Payam sich dem Richter direkt gegenüber setzte. Die beiden begannen ihr Gespräch auf Farsi. Ich beobachtete Payams Gesicht, das von Sekunde zu Sekunde besorgter wurde. Obwohl ich kein Wort verstand, kroch ein Unbehagen in mir hoch. Nach kaum zwei Minuten verließen wir wieder den Raum. Ihr wisst sicher, was jetzt kommt. Zahlen Sie 48.000 \$. Gehen Sie nicht über Los. Ziehen Sie keine 4.000\$ ein. Hier wurde die Bedeutung des zuvor gewählten Autowertes  von 16.000\$ schlagartig klar. Es ist im Iran gängige Praxis, die Strafe in einem Schmuggelfall auf das Dreifache des Warenwertes festzusetzen, weshalb auch die Sicherheitsleistung dieser Summe entsprach. Wir fuhren sprachlos zurück nach Norduz. Auf dem Weg kam mir in den Sinn, dass nun der Moment gekommen ist, mich bei der deutschen Botschaft zu melden. Ich wählte die Nummer und erklärte meine Situation. Die Frau am anderen Ende der Leitung sagte mir kurz und knapp, dass sie den Fall weitergeben werde und sich jemand bei mir melden würde. Kurz darauf meldete sich tatsächlich eine Frau Jäger bei mir. Sie wirkte am Telefon sehr eloquent und überaus rational, was in diesem Chaos eine seltsam beruhigende Wirkung auf mich hatte. Auch wenn sich aus diesem ersten Telefonat noch keine konkrete Lösung ergab, verblieben wir so, dass sie versuchen würde, ein offizielles Schreiben aufzusetzen. Vor allem sollte darin erklärt werden, dass ich aufgrund der kriegerischen Situation zwischen Pakistan und Indien gezwungen war, das Auto im Iran zu lassen. Das Problem dabei war jedoch, dass es zu diesem Zeitpunkt im Iran noch immer kein Internet gab. Sie war sich daher unsicher, ob sie mir den Brief überhaupt hätte zuschicken können. In Norduz angekommen wollten wir noch einmal das Gespräch mit Brille zu suchen, in der Hoffnung, unsere Position irgendwie zu verbessern. Doch aus diesem Plan wurde nichts. Stattdessen erzählte er uns von einer Abmachung, die der Zollchef wohl bereits auf kurzem Dienstweg mit dem Richter getroffen hatte. Man wollte die Angelegenheit schnell vom Tisch haben. Die Einigung sah eine Strafe von 55 \% des festgesetzten Autowertes vor. Das bedeutete im Umkehrschluss jedoch auch, dass ich offiziell dafür verurteilt würde, den Wagen geschmuggelt zu haben. Beiläufig erwähnte er zudem, dass er selbst die Anzeige erst sechs Tage vor meiner Ankunft an der Grenze gestellt hatte. Wäre ich nur eine Woche früher dort gewesen, wäre mir dieser ganze Spaß erspart geblieben. Wie das Schicksal manchmal so spielt. Ich war mein Leben lang stolz darauf, eigentlich nie Hass für jemanden empfunden zu haben. Meist kann ich es sehr gut nachvollziehen, warum Menschen handeln, wie sie handeln. In diesem Augenblick waren mir die Beweggründe aber vollkommen egal. Ich verspürte jetzt einen innigen Hass auf Brille. Mir ging nur durch den Kopf, wie dieser Mensch mit mir umgegangen ist: angefangen bei seinem „No, no problem“ am ersten Tag bis hin zu der Zusage, mich zum Gericht zu begleiten, was er dann doch nicht tat. Er wusste genau, dass ich dort ohne Persischkenntnisse nicht einmal erfahren werde, was mir vorgeworfen wird, und hat mich dennoch einfach allein mit dem Taxi dorthin fahren lassen. Ich überschlug den Schaden: 8.000€ Strafe, 5.000€ meine Kaution für das Carnet und das Auto, das ich nun nie wieder bekommen würde. Gesamthöhe des Schadens: 18.000 €, ungeachtet davon, dass mein Herz wirklich an diesem Mercedes hängt. Wir verließen das Grenzgelände und gingen ins Hotel.

\chapter*{Die bittere Träne der Freiheit}
\addcontentsline{toc}{chapter}{Die bittere Träne der Freiheit}

Nach der Rückkehr ins Hotel herrschte eine fast greifbare Anspannung. Payam war sofort in Telefonate vertieft, seine Stimme ein ständiges Murmeln im Hintergrund. Ich zog mich in das Nebenzimmer zurück, um meine Eltern zu erreichen. Es war eines dieser Gespräche, die man niemals führen möchte. Ich musste ihnen von der dramatischen Entwicklung berichten, von der astronomischen Kaution und der drohenden Realität, dass ich bald in einem iranischen Gefängnis landen könnte. Als ich das Gespräch beendete, trat ich ans Fenster. Mein Blick wanderte hinüber in Richtung Armenien. Dort drüben, nur ein paar hundert Meter entfernt, lag die Freiheit, die in diesem Moment so unerreichbar schien. In diesem Augenblick brach alles aus mir heraus. Ich weinte bitterlich, Tränen der puren Ohnmacht und Erschöpfung. Ich konnte mich nicht erinnern, wann ich das letzte Mal so die Kontrolle verloren hatte. Doch seltsamerweise fühlte es sich gut an. Als die Tränen versiegten, war die lähmende Trauer einer kühlen Entschlossenheit gewichen. Ich begriff, dass ich nicht einfach in der Ecke sitzen und auf mein Schicksal warten durfte. Ich wusste, ich musste irgendwie handeln, merkte aber auch, dass ich alleine nicht mehr Herr der Lage war. Obwohl ich noch nie ein großer Freund davon war, mich auf Social Media zu präsentieren, und meine Posts im Laufe der Reise eher seltener geworden waren, entschied ich mich nun doch, auf Instagram um Hilfe zu suchen. Ich veröffentlichte einen Beitrag, in dem ich nach jemandem suchte, der mir bei juristischen Problemen im Iran helfen konnte. Die Resonanz war gewaltig. Es meldeten sich unzählige Menschen und ich versuchte, jedem zu antworten, kam aber nach kurzer Zeit schlicht nicht mehr hinterher. Es ist wirklich spannend, wie Menschen in solchen Momenten reagieren. Die meisten fragten zwar nur besorgt, wie es mir geht, oder gaben den gut gemeinten, aber wenig hilfreichen Rat, mich bei der deutschen Botschaft zu melden. In einer solchen Extremsituation ist das natürlich nicht sonderlich förderlich. Doch zwischen all dem Rauschen gab es ein paar Nachrichten, die wirklich strukturiert waren und eine echte Option boten. In meiner Verzweiflung kontaktierte ich auch MC Basstard, einen deutschen Rapper und Iran-Aktivisten, mit dem ich zuvor schon ein paar Mal wegen meiner Reise Kontakt hatte. Er empfahl mir einen iranischen Anwalt aus Berlin. Das Telefonat mit ihm war jedoch ein Schlag in die Magengrube. Eines der ersten Dinge, die er mir sagte, war: „Ihnen muss klar sein, dass Sie jetzt eine Geisel sind.“ Auch wenn mein Verstand mir sagte, dass es rein rational noch keine Anzeichen für eine politische Geiselnahme gab, brennt sich so ein Satz tief ein, wenn man ohnehin schon mit dem Rücken zur Wand steht. Mir fiel wieder der Brief der Botschaft ein, auf den ich so dringend wartete. Da ich aber davon ausging, dass der Versand am fehlenden Internet scheitern würde, entschloss ich mich zu einem direkten Versuch beim Auswärtigen Amt in Berlin. Meine Hoffnung war simpel: Vielleicht konnten sie mir das Dokument von Deutschland aus direkt ausstellen. Nachdem ich mich durch die Warteschleife der Hotline gekämpft hatte, erreichte ich eine Mitarbeiterin und begann, meine Lage zu schildern. Doch kaum war das Wort „Iran“ gefallen, wurde ich unterbrochen. Sie erklärte mir in belehrendem Ton, dass für das Land eine Reisewarnung ausgesprochen sei. Ich bestätigte ruhig, dass mir das bewusst sei, und versuchte, in meiner Erzählung fortzufahren. Bevor ich jedoch zum Kern der Sache kommen konnte, fiel sie mir noch zwei weitere Male ins Wort, um mich erneut auf die Reisewarnung hinzuweisen. Schließlich platzte mir der Geduldsfaden; ich entgegnete ihr, dass ich das Prinzip verstanden habe, ich nun aber einmal in dieser Situation stecke und ihre Belehrungen in diesem Moment alles andere als hilfreich seien. Nachdem wir uns mühsam synchronisiert hatten, versprach sie, mich weiterzuleiten. Doch beim nächsten Kollegen bot sich das exakt gleiche Schauspiel. Er ließ mich nicht aussprechen und hielt mir insgesamt drei Mal die Reisewarnung vor. Mir ist völlig bewusst, dass ich keine großen Hilfeleistungen erwarten kann, wenn ich in ein Land einreise, für das eine offizielle Warnung besteht. Dennoch sollte es möglich sein, nach einer einmaligen Belehrung entweder zuzuhören, worin das eigentliche Problem besteht, oder direkt zu kommunizieren, wenn keinerlei Unterstützung geleistet werden kann. In einer derartigen Extremsituation ist dieses bürokratische Pingpong eine reine Tortur für jemanden, der in wirklich großen Schwierigkeiten steckt. Immerhin verlangte ich keine militärische Evakuierung, sondern lediglich einen kurzen Brief für das Gericht. Als ich ihm endlich begreiflich machen konnte, dass ich für das Gericht eine Bestätigung über den Indien-Pakistan-Konflikt benötigte, der mich zum Verbleib des Autos gezwungen hatte, kam nur ein trockenes: „Alles klar, ich gebe das weiter.“. Auf meine Nachfrage, wohin er das denn weitergibt, antwortete er lapidar: an die Botschaft. Die Botschaft, die ohnehin schon informiert war und mir wegen des Blackouts nichts schicken konnte. Das war das ernüchternde Resultat eines zwanzigminütigen Telefonats, das man auch in zwei Minuten hätte erledigt haben können. Den Rest des Abends verbrachte ich wie in Trance damit, Nachrichten zu schreiben, doch eine echte Lösung kristallisierte sich nicht heraus. Mein Kopf dröhnte von den stundenlangen Gesprächen; ich konnte keinen klaren Gedanken mehr fassen. Irgendwann schleppte ich mich rüber zu Payam. Wir redeten kurz über die weiteren Pläne. Er sagte, wir würden am nächsten Tag nach Täbris fahren, um Saleh zu treffen. Er würde schon eine Lösung finden. Ich dachte nur: „Okay“, ohne mir viel Hoffnung zu machen. Bevor ich wieder rüberging, erzählte Payam noch, dass die Amerikaner mit dem Flugzeugträger Abraham Lincoln im Persischen Golf angekommen seien. Ach ja, da war ja was... Ich legte mich ins Bett, unfähig, noch ein weiteres Wort zu wechseln. An Schlaf war nicht zu denken. Die Vorstellung, im Iran in den Knast zu kommen, war das eine. Damit würde ich irgendwie klarkommen. Ich wäre nicht der erste Ausländer, dem das passiert, und ich war mir sicher, dass die meisten Gefangenen dort normale Bürger sind, die mir eher helfen würden. Was aber viel schlimmer war als die bloße Zelle: Die Vorstellung, dass bald der Krieg ausbricht und ich dann als Gefangener direkt zum Spielball zwischen den Fronten werde. Dass ich dann nicht mehr nur wegen eines Autos feststecke, sondern als Faustpfand in einem militärischen Konflikt ende. Dieser Gedanke löste eine ganz andere Art von Angst aus. Erst nach Stunden der mentalen Erschöpfung schlief ich schließlich ein.

\chapter*{Das Gesetz der Schachtel}
\addcontentsline{toc}{chapter}{Das Gesetz der Schachtel}

Wir wachten am nächsten Morgen früh auf und machten uns auf den Weg nach Täbris. Da mein Pass weg war, gingen mir verschiedene Fluchtszenarien durch den Kopf. Das Auto und das Carnet waren mir in diesem Moment egal, es ging nur noch darum, hier irgendwie rauszukommen. Die Sicherheitslage hatte sich seit dem Vortag dramatisch verschlechtert. Ich war mir relativ sicher, dass im zentralen System keinerlei Vermerk über meine Gerichtsverhandlung existierte, da die iranischen Systeme untereinander nicht weit vernetzt sind. Mein Plan war daher, zur Ausländerpolizei zu gehen, den Verlust meines Passes vorzutäuschen und mit meinem zweiten deutschen Reisepass ein Exit Visum zu beantragen. Das war auch die Empfehlung von Frau Jäger aus der Botschaft. Ich wollte erst abwarten, was Saleh in Täbris erreichen konnte, aber das war definitiv mein Plan B, falls sich die Lage nicht schlagartig verbesserte. In Täbris angekommen holten wir Saleh ab. Er erklärte uns, dass wir direkt zum obersten Richter für Schmuggelangelegenheiten der Provinz fahren würden, den er durch seine Arbeit am Flughafen gut kannte. Das klang nach einem soliden Plan. Bevor wir das Gericht ansteuerten, hielten wir jedoch bei einer großen Konditorei. Inmitten von Bergen aus Gebäck und Torten besorgten wir für umgerechnet 10€ eine Schachtel der besten Pistazienkekse des Hauses. Ich war zwar verwundert über ein solches Mitbringsel, aber so schienen die Dinge hier eben zu laufen. Das Gericht war nur etwa 200 Meter entfernt. Am Eingang mussten wir unsere Handys abgeben, dann betraten wir das Büro des Richters. Er war ein sehr ernster Mann mit einer großen Narbe auf der Wange. An seiner linken Hand fehlten vier Finger, Verletzungen aus dem Krieg, die ihm vermutlich dabei geholfen hatten, in diese Position zu gelangen. Wir nennen ihn ab hier Scarface. Saleh erklärte ihm die Situation. Anfangs wirkte Scarface nicht besonders hilfsbereit und blickte nur finster drein. Erst als wir ihm ein Bild meines Mercedes zeigten und ich einige seiner Fragen auf Englisch beantwortete, taute er auf. Er begann zu lächeln, was ihn direkt sympathischer wirken ließ. Es gab Tee, die Stimmung lockerte sich und er begann zu telefonieren. Zuerst rief er den Richter in Jolfa an und sagte ihm, er solle den Fall entweder an den Zoll zurückgeben oder komplett schließen. Danach knöpfte er sich den Zollchef vor und erklärte ihm in deutlichem Ton, dass sie nicht jeden Schwachsinn bei ihm abladen sollten, immerhin seien im letzten Jahr 300 Fälle aus Norduz bei ihm gelandet. Zum Schluss telefonierte er sogar mir Brille und wies ihn an, mir meinen Pass zurückzugeben. Ich war fassungslos. Ohne Kontakte war ich gegen Mauern gerannt, aber mit Saleh an der Seite lief plötzlich alles glatt. Wir verließen das Gebäude und ich hatte ein breites Grinsen im Gesicht. Ich sah mich im Geist schon morgen mit dem Wagen aus dem Iran ausreisen. Es blieb jedoch das Problem mit meinem Visum, das genau an diesem Tag ablief. Saleh tätigte einen Anruf und wir fuhren zur Ausländerpolizei. Dort kannte Saleh offensichtlich jeden, es gab Umarmungen und herzliche Begrüßungen. Während oben die Details besprochen wurden, schickte Saleh Payam los, um Nachschub zu holen: neue Kekse für den Chef der Ausländerbehörde. Payam kam mit einer ähnlichen Schachtel zurück. Wir betraten das Büro, zeigten die Kopie meines Ausweises mit der handschriftlichen Notiz und erklärten die Lage. Der Chef war sehr freundlich, fragte, ob 20 Tage Verlängerung ausreichen würden und bestellte wieder Tee. Die drei unterhielten sich angeregt auf Farsi. Ich verstand kein Wort, aber das war mir egal, alle Probleme schienen sich in Luft aufzulösen. Dann kam ein Mitarbeiter mit ernster Miene herein und erklärte, dass sie das Visum ohne den Originalpass nicht im System verlängern könnten. Es gab eine kurze Diskussion, bis der Chef mit dem Finger auf den Tisch haute und klarmachte, dass eine funktionierende Lösung her müsse. Der Mitarbeiter brachte schließlich einen Zettel, den Saleh ausfüllte und ich unterschrieb. Man sagte mir, ich solle einfach mit meinem Pass wiederkommen, um alles nachzutragen, aber bis dahin sei alles in Ordnung. Beim Rausgehen wurde mir jedoch etwas klar. Jetzt stand im System, dass es ein Problem mit meinem Pass gab. Mein Plan B, mit dem Zweitpass einfach auszureisen, war damit schlagartig gefährlicher geworden. In diesem Moment verdrängte ich den Gedanken jedoch, schließlich wollte ich ja morgen mit meinem Mercedes das Land verlassen.

\chapter*{Zwischen Lammkoteletts und Existenzangst}
\addcontentsline{toc}{chapter}{Zwischen Lammkoteletts und Existenzangst}

Der nächste Morgen bei Payam bricht früh an. Er hat mir weiterhin sein Zimmer überlassen und schläft selbst auf der großen Couch im Wohnzimmer. Ohne Frühstück holen wir Saleh ab und fahren direkt zum Gericht, um die Angelegenheit endlich abzuschließen. Saleh geht allein zum Richter, während Payam und ich draußen im Flur warten. Nach einer Weile bringt ein Mitarbeiter Tee in das Richterzimmer, was ich für ein gutes Zeichen halte. Doch als Saleh schließlich herauskommt, dämpft er meine Erwartungen sofort. Er erklärt, dass der Richter zwar nicht persönlich gegen mich ist, den Fall aber nicht einfach fallen lassen kann. Da der Prozess bereits offiziell aufgenommen wurde, ist das Risiko für Korruptionsvorwürfe innerhalb des Systems zu hoch. Der Richter hat Saleh jedoch von einer neuen Möglichkeit erzählt: Der Zollchef in Täbris könnte das Verfahren stoppen, da der Zoll in diesem Fall der offizielle Kläger ist. Damit bleiben uns nur zwei Optionen. Entweder ich akzeptiere die Situation, riskiere ohne die 48.000 \$ die Haft oder wir versuchen, den Zollchef von einer Einstellung des Verfahrens zu überzeugen. Die Entscheidung fällt uns nicht schwer. Auf dem Weg zurück nach Täbris halten wir um 14 Uhr bei einem alten Bauern am Straßenrand und kaufen Äpfel. Es ist das erste Mal an diesem Tag, dass ich etwas esse. In Täbris angekommen, ändern sich die Pläne jedoch unerwartet. Saleh scheint den Zollchef nicht erreicht zu haben. Stattdessen steuern wir einen Neubau an, um die Schlüssel für seine neue Wohnung abzuholen. Wie üblich halten wir vorher bei einer Konditorei, um Süßigkeiten für die Übergabe zu besorgen. Wir betreten einen großen Versammlungsraum. Am Ende eines langen Tisches sitzt ein Immobilienmakler im schicken Anzug mit einem verschmitzten Lächeln. Neben ihm seine deutlich jüngere Assistentin, ganz in Schwarz gehüllt und sehr zurechtgemacht. Beide mustern mich interessiert, Ausländer sieht man hier wohl eher selten. Ich lächle höflich zurück, habe aber eigentlich keinen Kopf für Smalltalk. Saleh beginnt das Gespräch und erzählt ihnen schließlich meine ganze Geschichte inklusive des Gerichtsprozesses. Während sie plaudern, fangen sie an, ein wenig zu witzeln. „Bleib doch einfach hier und heirate“, sagen sie lachend. Die Assistentin kichert und erzählt von einer Freundin, die Englisch lernt, um einen Mann aus Europa zu heiraten. Eigentlich ist das eine normale Reaktion, aber in meiner Lage überkommt mich ein beklemmendes Gefühl. Ich fühle mich wie in einem schlechten Film. Ich stehe mit einem Bein im Gefängnis und hier werden Witze gemacht, von denen ich kaum ein Wort verstehe. Ein Dröhnen erfüllt meine Ohren, alles um mich herum wird verschwommen. Es ist unmöglich, in dieser absurden Situation einen klaren Gedanken zu fassen. Nach etwa 20 Sekunden lässt das Dröhnen nach. Die Assistentin sagt etwas, das Payam mir übersetzt: „So wie er guckt, ist er schon seit dem zweiten Weltkrieg hier eingesperrt.“ Ich merke, dass ich mich zusammenreißen muss, doch auch Saleh und Payam wirken nun besorgt. Die Leichtigkeit im Raum ist verschwunden. Wir erledigen die Formalitäten, nehmen die Schlüssel und besichtigen kurz die Wohnung. Danach fahren wir in ein gehobenes Restaurant und bestellen das teuerste Gericht auf der Karte: super zarte Lammkoteletts, die auf einer Art Schwert aufgespießt sind. Pro Person kostet das umgerechnet nur etwa 10 €. Während des Essens stabilisiert sich meine Verfassung etwas. Der enorme Stress der letzten anderthalb Wochen und der leere Magen hatten mich körperlich und mental komplett ans Limit getrieben. Mit jedem Bissen kehrt die Energie langsam zurück. Auf dem Weg, als wir Saleh nach Hause bringen, tippe ich eine Nachricht in meinen Übersetzer und zeige sie ihm: „Ich weiß, dass es absolut nicht selbstverständlich ist, was du für mich machst. Ich vergesse nie, wenn mir jemand einen Gefallen tut, vor allem nicht bei einer so großen Sache. Sag Bescheid, falls ich jemals etwas für dich tun kann, egal was es ist. Es tut mir leid, dass ich vorhin so die Fassung verloren habe. Es war einfach etwas zu viel.“ Saleh liest die Zeilen schweigend. Er sieht für einen Moment wirklich gerührt aus. An seinem Blick merke ich sofort, dass er es sehr zu schätzen weiß, dass ich ihm diese Anerkennung so direkt mitgeteilt habe. Um 22 Uhr treffen wir uns am Flughafen noch mit einem erfahrenen Zollinspektor. Er hilft uns dabei, ein Schreiben aufzusetzen, um den Zollchef davon zu überzeugen, das Verfahren endgültig einzustellen. Wir verbringen etwa zwei Stunden in seinem kleinen Büro, während er das Schriftstück verfasst. Immer wieder stellt er Detailfragen, die Payam geduldig für ihn beantwortet. Als wir das Büro schließlich verlassen, sind wir die einzigen verbliebenen Gäste im Terminal. In der Leere des Gebäudes fällt mir auf, dass uns ein Mann vom iranischen Geheimdienst halb unauffällig folgt. Ich denke mir in diesem Moment nicht viel dabei; ich bin einfach nur froh, dass dieser lange Tag zu Ende geht.

\chapter*{Bürokratie im Leerlauf}
\addcontentsline{toc}{chapter}{Bürokratie im Leerlauf}

Der Donnerstag markiert im Iran den Beginn des Wochenendes. Das bedeutete für mich vor allem: Stillstand. Wenn die Behörden schließen, passiert in diesem Land absolut nichts mehr. Es war ein seltsames Gefühl. Einerseits nagte es an mir, dass zwei Tage lang nichts vorangehen würde, während meine Zeit und mein Visum knapp wurden. Andererseits war es dringend nötig, einmal auszuschlafen und nach dem Stress der letzten anderthalb Wochen wieder zu Kräften zu kommen. Wir nutzten die Zeit für ein paar organisatorische Dinge. Das Schreiben, das der Zollinspektor am Vorabend aufgesetzt hatte, war nur handschriftlich verfasst. Im Iran ist es eine gängige Kombination, dass man erst ein handgeschriebenes Dokument erstellt und dann einen Copy Shop aufsucht, um es abtippen zu lassen. Also machten wir uns auf die Suche nach einem Laden, der diesen Service anbietet. Bei dieser Gelegenheit druckten wir auch direkt eine Kopie meines Carnet de Passage und das Schreiben der deutschen Botschaft aus, das ich am Dienstagmorgen erhalten hatte. Den Rest des Tages verbrachten wir auf dem Bazar, um noch einige Erledigungen zu machen. Der Freitag verlief ruhig. Wir trafen uns lediglich mit Saleh und dem Zollverantwortlichen des Flughafens, um die gesammelten Papiere noch einmal final durchzugehen. Wir wollten sicherstellen, dass alles bereit ist, wenn wir am nächsten Tag zum Zoll in Täbris gehen. Die Einschätzung der Experten vor Ort war positiv: Die Unterlagen, die wir nun zusammen hatten, sollten ausreichen, damit die Anzeige gegen mich fallen gelassen wird. Am nächsten Morgen machten wir uns auf den Weg zum Zoll, doch die Ernüchterung folgte sofort: Der Chef war nicht in seinem Büro. Man sagte uns, er würde erst am Nachmittag wiederkommen. Um die Zeit sinnvoll zu nutzen, fuhren wir noch einmal bei Scarface vorbei, um den aktuellen Stand aus seiner Sicht zu erfahren. Sein Urteil war jedoch wenig ermutigend. Er teilte uns mit, dass ich für meine letzte Verteidigung erneut zum Gericht gehen müsse. An der Kaution gebe es nichts zu rütteln; sie sei nicht verhandelbar. Wenn ich die Summe nicht aufbringe, bliebe der Gang ins Gefängnis die einzige Konsequenz. Auf dem Rückweg vom Zollamt vibrierte mein Handy. Eine Sprachnachricht von Reza, einem Freund aus Hamburg, mit dem ich vor meiner Reise zusammengewohnt habe. Er hatte einen befreundeten Mullah in Teheran gebeten, mir aus meiner Lage zu helfen. Das Ergebnis war jedoch negativ: Der Kontakt lehnte jede Unterstützung ab. Jemand aus dem Team des Mullahs hatte mich im System geprüft und dabei festgestellt, dass eine sogenannte Red Flag auf mich gesetzt wurde. Das bedeutet, dass der Geheimdienst ein Auge auf mich hat. Wer diesen Vermerk veranlasst hat, konnten sie nicht sagen, da bereits die Abfrage dieser Information im System nachverfolgbare Spuren hinterlässt. Reza riet mir eindringlich, mich bedeckt zu halten. Er spricht aus Erfahrung: Als Deutsch-Iraner saß er selbst vor einigen Jahren für neun Monate im iranischen Gefängnis, weil man ihn der Spionage verdächtigte. Es war einer dieser Momente, in denen man denkt, es könne nicht mehr weiter bergab gehen, und dann kommt die nächste Hiobsbotschaft um die Ecke. Doch nach den letzten anderthalb Wochen fühlte sich diese Warnung fast wie eine weitere Kleinigkeit an, die ich nun auch noch auf meinen Schultern trug. Mein rationaler Gedanke war: Hätte der Geheimdienst wirklich etwas Handfestes gegen mich, hätten sie mich wahrscheinlich längst geholt. Dennoch hatte ich absolut keine Lust auf ein stundenlanges Verhör bei meiner geplanten Ausreise. Ich entschied mich dazu, niemandem davon zu erzählen. Obwohl Payam und ich mittlerweile gute Freunde waren und offen über alles sprachen, war die Gefahr zu groß, ihn einzuweihen. Hätte er aus Angst den Kontakt abgebrochen, wäre ich wieder völlig allein gewesen, ohne jeden Ansatzpunkt zur Lösung meiner Situation. Auch meinen Eltern und Freunden in Deutschland erzählte ich nichts, um sie nicht noch mehr zu beunruhigen. So fuhren wir schweigend zurück zum Zoll, in der Hoffnung, den Chef am Nachmittag endlich anzutreffen. Doch auch dieses Mal war sein Platz leer. Es ist eine frustrierende Erfahrung, wie häufig im Iran absolute Schlüsselpersonen einfach nicht an ihrem Arbeitsplatz anzutreffen sind, wenn es darauf ankommt.

\chapter*{Ein Lichtblick mit Hindernissen}
\addcontentsline{toc}{chapter}{Ein Lichtblick mit Hindernissen}

Der nächste Tag begann vielversprechend. Wir fuhren frühmorgens direkt zum Zoll und hatten diesmal tatsächlich Erfolg: Der Assistent des Chefs nahm sich Zeit, uns anzuhören. Nach dem üblichen Hin und Her zeichnete sich ein echter Fortschritt ab. Der Chef entschied, eine offizielle Stellungnahme von der Grenze in Norduz anzufordern. Für mich war das ein kleiner Lichtblick in diesem bürokratischen Dickicht, auch wenn wir erst am Dienstag wiederkommen sollten. Da wir keine Zeit verlieren wollten, holten wir Saleh ab und machten uns sofort auf den Weg zurück nach Norduz. Wir wollten persönlich sicherstellen, dass die Kommunikation zwischen den Behörden nicht im Sande verläuft. An der Grenze angekommen, suchten wir direkt das Gespräch mit Brille. Doch von Kooperation war keine Spur. Obwohl Scarface ihn angewiesen hatte, mir meinen Pass zurückzugeben, flüchtete er sich in eine dreiste Notlüge und behauptete, alle meine Dokumente lägen bereits beim Gericht. Wie sich später herausstellte, war das eine glatte Lüge, die man ihm in diesem Moment schon an seinem Verhalten ansah. Ich verstand mal wieder kein Wort der hitzigen Diskussion, merkte aber, wie die Stimmung gefährlich kippte. Saleh redete immer energischer auf ihn ein, während Brille den Moment sichtlich genoss. Er schien sich an der Macht zu berauschen, die er in diesem Augenblick über mich hatte. Payam und Saleh schwiegen sich danach darüber aus, was Brille genau gesagt hatte. Warum sie mir die Details vorenthielten, bleibt mir bis heute ein Rätsel. Ich wusste in diesem Moment einfach nicht, wo das eigentliche Problem lag. Heute nehme ich an, dass uns ein entscheidendes Puzzleteil fehlte: Der Zoll in Täbris konnte mir faktisch gar nicht helfen, da die wahre Entscheidungsgewalt in Teheran lag. Doch zu diesem Zeitpunkt hatten wir dorthin keinerlei Verbindung und versuchten verzweifelt, die Dinge vor Ort zu regeln. Im Anschluss unterhielt sich Saleh relativ lange mit dem schnurrbärtigen Zollchef von Norduz. Am Ende wirkten beide fast zufrieden mit dem Gespräch. Der Chef gab Saleh sein Wort, dass sie nicht in Berufung gehen würden. Für mich änderte das jedoch wenig; nach all den Lügen war mir klar, dass ich diesen Menschen und ihren Versprechen in keiner Weise vertrauen konnte. Nach diesem kräftezehrenden Tag traten wir die fast dreistündige Rückreise an. Es war ein erschöpfender Weg, und tief im Inneren ahnte ich bereits, dass es nicht das letzte Mal gewesen sein würde, dass wir diese Strecke hinter uns bringen mussten. Trotz der Ungewissheit stand unser Entschluss fest: Wir würden am nächsten Tag wieder beim Zoll in Täbris erscheinen. In unserer Lage gab es ohnehin nichts mehr zu verlieren.

\chapter*{Ein Anwalt muss her}
\addcontentsline{toc}{chapter}{Ein Anwalt muss her}

Am nächsten Morgen fuhren wir wie geplant zum Zoll in Täbris. Wie bereits erwartet, wurde uns vom Assistenten gesagt, dass die Anzeige nicht zurückgenommen werden kann. Er las Payam auch die angeforderte Einschätzung aus Norduz vor. Brille hatte darin nun deutlich geschrieben, dass er auch nicht mehr davon ausgehe, dass ich das Auto schmuggeln wollte. Auch räumte er ein, dass der zuvor berechnete Wert von 16.000 \$ deutlich zu hoch angesetzt war. Dennoch könne man die Anzeige nicht zurückziehen, da dem Zoll nicht bekannt sei, wofür das Auto in den letzten Monaten meiner Abwesenheit verwendet wurde. Man versicherte uns aber: Falls das Gericht nachfragen sollte, würde man definitiv ein gutes Wort für mich einlegen. Auch wenn die Lage nicht wirklich besser geworden war, erfuhren wir, dass wir eine Bescheinigung darüber brauchten, wo der Wagen die ganze Zeit verblieben war. Wir riefen Razool in Schiras an, bei dem der Mercedes geparkt war. Er erzählte von einem guten Kontakt zu einem Mr. Haghidat vom Automobil Club. Dieser könne mir eine Bestätigung mit offiziellem Stempel ausstellen. Payam rief ihn an. Dieser Mann war wirklich ein Segen: Er machte direkt am nächsten Tag alles fertig und schickte es uns per E-Mail zu. Nach all unseren Versuchen gab es keinen Ausweg mehr: Ich musste vor Gericht. Dafür brauchte ich einen Anwalt. Wir hatten uns bereits über verschiedene Möglichkeiten informiert. Die meisten Anwälte machten uns klar, dass ich in einem Riesenproblem steckte. Sie forderten ein Honorar von mindestens 10 \% des Streitwerts, was in meinem Fall zu diesem Zeitpunkt 4.800 \$ entsprach. Ich war nicht bereit, eine solche Summe zu zahlen. Bereits ein paar Tage vorher sendete mir ein alter Bekannter, den ich in Portugal kennengelernt habe, eine Nummer von einem Anwalt aus Jolfa. Er hatte diesen nach meinem Instagram-Post über mehrere Ecken empfohlen bekommen. Payam rief ihn an. Dieser sagte einfach nur: „Kommt rum und wir gucken mal.“ In Jolfa angekommen fanden wir nach ein paar Orientierungsproblemen schließlich die richtige Adresse. Der Anwalt empfing uns direkt draußen auf der Straße. Er war ein nett lächelnder Mann, der auf den ersten Blick eher zurückhaltend wirkte. Wir gingen in sein Büro, übergaben ihm alle Dokumente und Payam erzählte ihm die gesamte Geschichte. Am Ende machte Payam ihm unmissverständlich klar, dass von mir finanziell nicht viel zu holen sei. Wir verabschiedeten uns und fuhren anschließend zurück in das bereits wohlbekannte Hotel in Norduz.

\chapter*{Das Siegel aus Tinte und Tränen}
\addcontentsline{toc}{chapter}{Das Siegel aus Tinte und Tränen}

Am nächsten Morgen trafen wir uns an einem großen Kreisverkehr in Jolfa. Der Anwalt sah heute deutlich herausgeputzter aus als noch am Vortag. Er erklärte mir, dass meine Situation juristisch eigentlich ziemlich eindeutig sei und die Indizien auf meiner Seite stünden. Sein Angebot für die Verteidigung lag bei gerade einmal 200 Euro. Im Vergleich zu den 4.800 Dollar, die wir nach den ersten Gesprächen mit anderen Anwälten befürchtet hatten, war das ein gewaltiger Unterschied. Wir sagten sofort zu. Er holte eine Vollmacht hervor, die ich unterschrieb und mit meinem Fingerabdruck besiegelte. Wir machten ihm noch einmal unmissverständlich klar, dass unser größtes Problem die Kaution sei, weil ich unter keinen Umständen ins Gefängnis gehen wollte. Er versprach, zum Gericht zu gehen und das weitere Vorgehen direkt mit dem Richter zu klären. Da wir an diesem Tag nicht viel mehr ausrichten konnten, fuhren wir am Aras entlang, dem Grenzfluss zu Armenien. Wir kamen in ein kleines Dorf, das fast nur aus Fischern zu bestehen schien. Payam kannte dort einige Leute, weil er früher häufig Geschäfte zwischen den lokalen Fischern und internationalen Abnehmern vermittelt hatte. Die Qualität des Fangs war beeindruckend. Wir unterhielten uns darüber, ob man diesen Fisch womöglich nach Russland oder Deutschland exportieren könnte. Am Ende bestellten wir zweimal gebratenen Fisch vom heutigen Fang. Es war einfach köstlich. Ich liebe diese Momente in entlegenen Dörfern, in denen man die Einfachheit einer Küche mit frischen lokalen Zutaten genießen kann. Doch genau dort erreichte uns der Anruf des Anwalts. Er hatte bereits eine erste Rückmeldung vom Gericht: Der Richter hatte uns einen offiziellen Termin für Montagmorgen gegeben. Damit war das Ende der Ungewissheit datiert, was die Stimmung am Tisch sofort veränderte.

Auf dem Rückweg nach Täbris fingen meine Gedanken jedoch an, verrückt zu spielen. Die Konfrontation mit dem Gericht in Jolfa war nun unausweichlich geworden. Die Lage fühlte sich bedrückend an:

\begin{itemize}
  \item Ich musste weiterhin damit rechnen, ins Gefängnis zu gehen, da ich keine Möglichkeit sah, die Kaution von 48.000 Dollar aufzubringen.
  \item Der Flugzeugträger Abraham Lincoln lag weiterhin direkt vor dem Persischen Golf.
  \item Payam versuchte zwar, Leute anzurufen, um Optionen für die Kaution zu finden, aber nichts davon klang in irgendeiner Form hilfreich.
\end{itemize}

Ich wusste nicht mehr weiter und erinnerte mich an Zahra. Sie ist eine georgisch iranische Tierärztin aus dem Iran, die jetzt seit einigen Jahren in Tiflis lebt. Wir hatten eine Romanze, auf die ich nicht genauer eingehen möchte. Wir waren zwar nicht im Streit, aber gefühlt auch nicht im Guten auseinandergegangen. Mir fiel ein, dass sie häufiger ihre leerstehende Eigentumswohnung in Teheran erwähnt hatte. Der bloße Gedanke, ausgerechnet sie um Hilfe bitten zu müssen, kratzte massiv an meinem Ego. Normalerweise ziehe ich mich nach bewährter Münchhausen-Manier eigenhändig aus dem Sumpf. Aber diesmal fühlte es sich an, als hätte man mir den Kopf rasiert. Die Entscheidung, Zahra in meiner Verzweiflung um Hilfe zu bitten, löste etwas in mir aus. In dem Moment, als ich den Entschluss im Kopf fällte, liefen mir plötzlich die Tränen über das Gesicht. Ich konnte nichts dagegen tun. Es war kein richtiges Weinen, die Tränen strömten einfach nur aus meinen Augen. Ich hatte so etwas vorher noch nie erlebt. Ich denke, Payam bemerkte es am Steuer, aber wir schwiegen eine Zeit lang. Als er mich schließlich etwas fragte, war ich überrascht, wie klar ich antworten konnte, obwohl meine Augen noch tränten. Nach einiger Zeit hielten wir an einer Tankstelle und holten uns etwas Nervennahrung in Form von süßem Gebäck. Ich fühlte mich danach etwas besser. Es ist erstaunlich, wie Weinen auf die menschliche Psyche wirken kann: Als wir etwa eine Stunde später bei Payam in Täbris ankamen, fühlte ich mich richtig gut, deutlich besser als die Tage zuvor. Wir aßen zu Abend und ich legte mich nach dem ganzen Stress etwas hin. Nach ein paar Stunden ging ich wieder ins Wohnzimmer. Payam berichtete mir, dass der Anwalt sich gemeldet hatte. Er hatte es tatsächlich vollbracht: Der Richter hatte entschieden, die Kaution auszusetzen. Vorerst zumindest, wie es sich später noch rausstellen sollte. Durch diese Nachricht erübrigte es sich schließlich auch, Zahra tatsächlich zu kontaktieren.

\chapter*{Warten, warten}
\addcontentsline{toc}{chapter}{Warten, warten}

Mittwoch. Ein Feiertag vor dem Wochenende und der Beginn einer weiteren Phase des Wartens. Vor Samstag können wir rein gar nichts tun. Es ist genau das gleiche üble Gefühl wie in der Woche zuvor: eine Mischung aus Ohnmacht und dieser nervösen Unruhe, die entsteht, wenn man zur Untätigkeit verdammt ist. Um die Zeit irgendwie zu füllen und den Kopf freizubekommen, machen wir ein paar Ausflüge. Einer davon führt uns nach Kandovan, in das iranische Kapadokien. Die uralten, in den Fels gehauenen Behausungen sind zwar beeindruckend, aber im Hinterkopf tickt ständig die Uhr, die eigentlich stillsteht. Am Samstag stehen wir erneut beim Zoll-Assistenten in Täbris. Wir erhoffen uns eine schriftliche Stellungnahme für das Gericht, ein Dokument, das meine Situation entlasten könnte. Doch die Antwort bleibt trocken und bürokratisch: Uns wird das Schreiben verwehrt. Es gibt keinen Fortschritt, nur die Erkenntnis, dass wir uns weiter im Kreis drehen. Also heißt es wieder: warten. Am Sonntag, während wir gerade ein paar Einkäufe erledigen, klingelt das Handy. Mein Anwalt ist dran. Die Nachricht ist kurz und ernüchternd: Die Verhandlung wurde erneut verlegt, diesmal auf Dienstag um 8 Uhr morgens. Geil, noch mehr Warten. Ich kann mich vor Freude kaum halten. Ehrlich gesagt kann ich mich nicht mehr genau daran erinnern, was danach bis zum Gerichtstermin passiert ist. Ich bin in eine leicht depressive Phase verfallen. Die ständige Ungewissheit und das Gefühl, im System festzustecken, haben mich mürbe gemacht.

\chapter*{Tag des Gerichts}
\addcontentsline{toc}{chapter}{Tag des Gerichts}

Es war so weit. Mein Termin beim Richter stand endlich an. Etwas nervös betrat ich das Richterzimmer, das ich bereits von meinen vorherigen Besuchen kannte. Die Szenerie war unverändert: Der Richter thronte hinter seinem Holztisch, während die Metallstühle drumherum noch immer sorgfältig in Plastikfolie eingeschlagen waren. Ich nahm den Platz direkt gegenüber vom Richter ein. Neben mir saß Mr. Translation. Payam hatte sich zunächst rechts in die Stuhlreihe gesetzt, doch der Richter bat ihn kurz darauf, das Zimmer zu verlassen. Zwei Schlüsselfiguren fehlten zu Beginn noch: Mein Anwalt und Brille. Wir fingen also ohne sie mit der Aufnahme meiner Personalien an. Dass das im Iran kein rein bürokratischer Akt ist, merkte ich spätestens bei der Religionsfrage. "What is your religion?", fragte Mr. Translation. "Agnostic", antwortete ich. "What's this?", erwiderte er mit einem völlig verwirrten Blick. In diesem Moment erschien es mir wenig zielführend, eine philosophische Abhandlung über das Konzept des Agnostizismus zu halten. Also sagte ich bloß: "Atheist." "You can't say this here!", antwortete er mit aufgerissenen Augen. Okay, dann bin ich eben: "Christ." Da ich sogar zwei Mal getauft bin, entsprach das ja gewissermaßen der Wahrheit. Im Nachhinein erfuhr ich, dass es für dieses Protokoll nur vier passable Antworten gibt: Schiit, Sunnit, Christ oder Jude. Kurz darauf betrat der Anwalt den Raum und begann sofort, meine Gerichtsunterlagen zu studieren. Wirklich gut vorbereitet wirkte er nicht, aber für das Geld konnte ich wohl auch nicht mehr erwarten. Der Richter klärte mich darüber auf, dass ich das Recht habe, zu schweigen und meinen Anwalt für mich sprechen zu lassen. Ich verneinte und fing an, die gesamte Geschichte von Anfang an zu erzählen. Während ich sprach, nahm mein Anwalt einen DIN A4 Notizblock und begann mitzuschreiben. Ich berichtete, wie ich am 17. März in den Iran gekommen war. Ich erzählte, wo ich mich die ersten zwei Monate aufgehalten hatte. Dass ich eigentlich nach Nepal wollte und wie der Indien-Pakistan-Konflikt mir den Weg versperrte. Auch die Bekanntschaft mit Alen, dem Kroaten, erwähnte ich. In der Zwischenzeit betrat auch Brille den Raum. Er trug meine Akte bei sich, in der mein Reisepass fest eingeklebt war. Von wegen, das Dokument liege bereits beim Gericht, wie er mir noch vor Tagen weismachen wollte. Ich ließ mich von seiner Anwesenheit nicht ablenken und erzählte weiter. Für den Richter war ein zentraler Punkt, wo das Auto abgeblieben war. Er wollte wissen, ob es auf einem Privatgrundstück stand , ob jemand anderes Zugang dazu hatte und ob in meiner Abwesenheit Reparaturen durchgeführt wurden. Ich erwähnte auch, dass in meinem Carnet de Passage bei der Einreise keinerlei Ausreisedatum notiert worden war. Das war ein bedeutender Punkt, denn gesetzlich wären die Grenzbeamten dazu verpflichtet gewesen, mir einen genauen Zeitpunkt für die Ausreise des Autos zu nennen. Der Richter verlas den dazugehörigen Gesetzestext. Brille gefiel dieser Moment ganz und gar nicht. Nachdem der Richter fertig war, verließ er den Raum, um zu telefonieren. Ich weiß nicht genau, mit wem er sprach, jedenfalls kam er kurz darauf wieder rein und versuchte, dagegen zu argumentieren. Es wirkte nicht so, als ob das beim Richter auf Wohlwollen stieß. Schließlich behauptete er einfach dreist, auf meinem Carnet stehe, dass das Auto nur einen Monat im Land bleiben dürfe. Als ich ihm die Kopie vorhielt, konnte er die Stelle scheinbar nicht finden. Ich beendete meine Verteidigung mit der Aussage, dass ich bloß ein Reisender sei, der das iranische Gesetz einfach nicht besser kannte und auf die falschen Leute gehört hatte. Jetzt war Brille an der Reihe. Als Ankläger versuchte er etwa zehn Minuten lang, den Richter davon zu überzeugen, dass ich wegen Schmuggels bestraft werden müsse. Der Richter schien sich nicht umstimmen zu lassen. Sichtlich enttäuscht setzte Brille sich wieder auf seinen Platz und fing an, in denselben Notizblock zu schreiben, in dem der Anwalt meine Aussagen protokolliert hatte. Bevor ich schließlich die Unterschriften unter das Protokoll setzte, wies uns der Richter an, dass er den originalen Brief bezüglich des Autoverbleibes haben wolle, da dieser bisher nur ausgedruckt vorlag. Außerdem verlangte er eine Quittung über die Reparaturen  und eine Übersetzung meines Carnets bis Samstag. Er versicherte uns, dass dann alle nötigen Beweise beisammen wären, um zu belegen, dass ich den Wagen nicht geschmuggelt habe und somit unschuldig sei. Der Richter erklärte nun, wie es weitergeht, doch Mr. Translation übersetzte davon absolut gar nichts. Allgemein habe ich im Verlauf des Verfahrens nur die absolut notwendigsten Punkte mitgeteilt bekommen. Als Brille mit seiner Niederschrift fertig war, wurde ich nach vorne gerufen. Ich sollte nun alle Aussagen, die mein Anwalt aufgeschrieben hatte, unterschreiben und mit meinem Fingerabdruck stempeln. Nach circa zwei Stunden verließen wir das Richterzimmer. Ich fühlte mich gut. Es war wie eine Befreiung. Die Session war scheinbar für alle ziemlich anstrengend gewesen, daher gingen wir zusammen nach draußen, ohne viel zu reden. Der Anwalt sagte noch ein paar Dinge zu Payam und verabschiedete sich. Da ich keinerlei Informationen über das weitere Vorgehen übersetzt bekommen hatte, ging ich fest davon aus, dass ich das Land in ein paar Tagen endgültig mit meinem Auto verlassen könnte. Es schien ja alles ziemlich gut gelaufen zu sein. Als der Übersetzer schon gehen wollte, sprach ich ihn noch einmal darauf an. Mr. Translation erklärte mir mit einem leichten Grinsen, dass es mindestens einen Monat dauern würde. Sein Grinsen wurde deutlich breiter, als er hinzufügte, dass es wahrscheinlich doch eher vier bis fünf Monate werden. Während er das in seiner seltsam freudigen Art sagte, hatte ich den kurzen Impuls, ihm ins Gesicht zu schlagen. Was war verkehrt mit diesem Typen? Es ist sein Job, die Aussagen im Gericht für mich zu übersetzen, er kommt dem nicht nach und macht sich am Ende noch über meine Situation lustig? Die Zeitangaben kamen wie folgt zustande: Nachdem das Urteil geschrieben ist, hat der Zoll 20 Tage Zeit, dagegen in Berufung zu gehen. Im Idealfall dauert es also etwa einen Monat. Der Richter erwähnte aber auch, dass der Zoll üblicherweise Berufung einlegt. Das kann dann gut und gerne ein halbes Jahr dauern. Und dann war da noch eine Sache. Die Kaution war nur bis zum Urteil ausgesetzt. Das erzählte man mir auch erst jetzt. Das heißt, sobald das Urteil steht und die 20 Tage Frist beginnen, müsste ich im Gefängnis warten. Falls der Zoll das Urteil anfechtet, säße ich die gesamte Zeit über im Knast. Eine positive Wendung gab es immerhin: Die Kaution wurde vom dreifachen auf den einfachen Fahrzeugwert gesenkt. Da die vom Krieg bedrohte Landeswährung weiter gefallen war, entsprach das nun etwa 12.000 Euro. Zwar keine 48.000 Dollar mehr, aber dennoch keine Summe, die ich einfach so rumliegen habe. Die gesamte Euphorie nach dem erfolgreichen Gerichtstermin war zerschlagen. Der Loop in meinem Kopf sprang wieder an: Nach einer positiven Sache folgt im Iran direkt die nächste Hiobsbotschaft.

\chapter*{In Khoy, um Khoy und um Khoy herum}
\addcontentsline{toc}{chapter}{In Khoy, um Khoy und um Khoy herum}

Auf dem Weg nach Täbris besprachen wir das weitere Vorgehen. Wir waren fast am Ziel, aber die Kaution stand noch aus. Wir hatten vorher schon darüber gesprochen, dass es im Iran üblich ist, dass jemand anderes die Kaution aufbringt, wenn man sie selbst nicht zahlen kann. Allerdings kassieren diese Leute üblicherweise 15 bis 20 \% des Betrages. Da die Kaution an sich herabgestuft worden war, wurde dieser Ausweg plötzlich erschwinglich. Payam rief seinen Freund Nasir aus der Stadt Khoy an. Nasir hat eine größere Fläche Land, die mindestens das Dreifache meiner Kaution wert ist. Sie einigten sich darauf, die Sache für 10 \% über die Bühne zu bringen. Angesichts der Wahl, entweder 1.200 Euro zu zahlen oder einen Krieg im Gefängnis miterleben zu müssen, fiel mir die Entscheidung denkbar leicht. Wir einigten uns also darauf, am Samstag gemeinsam zum Gericht zu fahren, um alles zu regeln. Wieso Samstag? Weil am Mittwoch natürlich wieder ein Feiertag war. Also mal wieder warten....

Es ist wirklich spannend zu beobachten, wie sich das Gemüt im Laufe einer solchen Situation verändert. Am Anfang ist man voll fokussiert, man will jeden einzelnen Schritt aufnehmen und verstehen. Man hat das Gefühl, dass man in diesem Moment wirklich alles tun muss und irgendwie auch etwas tun kann, um die Situation zu lösen. Dieser Fokus nimmt mit der Zeit ab. Fast vier Wochen nachdem ich an die Grenze kam, blieb nicht mehr viel davon übrig. Ich hatte mich wirklich daran gewöhnt, Dinge einfach hinzunehmen. Versteht mich nicht falsch: Es ist kein Gefühl des Aufgebens, sondern eher die Erkenntnis, dass man nicht wirklich etwas bewegen kann. Man verfällt in eine Ist-Mir-Egal-Haltung. Was passiert, das passiert. Man ist ausgelaugt und hat gelernt, dass Hoffnung nichts bringt. Frei nach dem Motto: ein Schritt vor, zwei zurück.

Es war so weit, Samstag. Mal wieder. Wir fuhren gut gelaunt früh am Morgen nach Khoy, um Nasir abzuholen. In Khoy stieg ein relativ kleiner, freundlich lächelnder Mann ein, dem sein Bart bis kurz unter die Augen wuchs. Er hatte einen Wrap dabei, der traditionell mit Kartoffeln, Ei und Kräutern befüllt war, und bot uns ein Stück an. Wir verneinten höflich und machten uns auf den Weg zum etwa 80 Kilometer entfernten Gericht in Jolfa. Dort angekommen, erklärte uns der geschniegelte Assistent, dass wir zum Tazirat (Schmugglergericht) in Khoy müssen. Da Nasirs Land in einer anderen Region liegt, könne man es hier nicht als Kaution annehmen. Vorher müssten sie jedoch eine Anfrage stellen, und das Ganze müsse innerhalb eines Tages abgewickelt werden, sonst drohe der Knast. Solch eine Aussage schockierte mich mittlerweile kaum noch. Wir sollten am nächsten Morgen wiederkommen, um die Anfrage für Khoy abzuholen. Wir brachten Nasir zurück zu seinem Auto und fuhren ins wohlbekannte Hotel in Norduz an der Grenze. Am nächsten Tag fuhren wir früh morgens zum Tazirat in Jolfa und dann weiter nach Khoy, wo Nasir bereits auf uns wartete. Dort gaben sie uns einen Brief für das Land-Department. Zurück beim Tazirat sagte man uns jedoch, dass das Land so nicht akzeptiert werden kann, weil es Ackerland und kein Bauland sei. Payam machte Witze darüber, dass es wohl wieder auf Samstag hinausläuft. Ich machte ihm klar, dass ich für solche Späße gerade wirklich nicht in der Stimmung war. Wir riefen Saleh an, dieser rief Scarface an, und dieser sprach mit meinem Richter. Wir kontaktierten den Anwalt, damit dieser ins Tazirat in Jolfa ginge, um der Sache nachzugehen. Er sagte, er sei eh in der Nähe und mache das direkt. Alles dauerte wieder eine gefühlte Ewigkeit. Ich musste wieder an Zahra denken, verwarf den Gedanken vorerst aber wieder. Der Anwalt rief zurück und sagte, der Assistent bleibe dabei: Bei solch einer Summe könne Ackerland nicht akzeptiert werden. Na toll... Wir riefen erneut den Anwalt an und sagten ihm, er solle persönlich mit dem Richter sprechen. Er willigte ein. Erneut verging eine Ewigkeit, bis der Rückruf kam: Das Ackerland wird als Kaution doch akzeptiert. Um 13:55 Uhr bekamen wir den neuen Brief fürs Land-Department. Da um 14 Uhr sämtliche staatlichen Einrichtungen schließen, konnten wir an diesem Tag nichts mehr erreichen. Wir machten uns auf den Weg zu Nasirs eigentlicher Einkommensquelle. Er betreibt eine Art Badehaus in den Bergen kurz vor der türkischen Grenze. Dort wird Wasser aus heißen Quellen in große Becken hochgepumpt. Natürlich gibt es separate Männer- und Damen-Bäder. Wir entspannten uns einige Zeit im heißen Wasser und zogen uns danach in ein kleines Häuschen zurück, das Nasir uns dort zur Verfügung gestellt hatte. Neben einem Mitarbeiter waren Payam und ich die einzigen, die die Nacht dort verbrachten, weil es nachts in den Bergen einfach noch sehr kalt war. Ich platzierte meine dünne Matte direkt neben dem Heizstrahler, auf dessen anderer Seite Payam schlief. Anders wäre es eine wirklich ungemütlich kalte Nacht geworden. Der nächste Morgen brach an. Wir fuhren durch die sonnenbestrahlten kurdischen Berge zurück nach Khoy zum Land-Department. Dort wurde uns mitgeteilt, dass auf dem Land noch eine Kaution offen sei. Nasir erzählte uns, dass er diese mal für seinen Bruder gestellt hatte und sie eigentlich bereits vor acht Jahren getilgt worden war; das schien aber nicht überall im System angekommen zu sein. Wir fuhren zum Hauptgericht in Khoy, um herauszufinden, wo das Problem lag. Da die nächsten regulären Termine erst in einer Woche frei waren, holten wir eine Schachtel Süßigkeiten und fuhren zu einem Büro, das Termine ausstellen kann. Dank der Süßigkeiten bekam Nasir einen direkten Termin am gleichen Tag. Wir warteten zwei Stunden am Hauptgericht und wurden dann zum Tazirat geschickt. Wir mussten noch bei der Bank vorbei, die ebenfalls in die Thematik involviert war. Um 14:10 Uhr kamen wir am Tazirat an. Es war zwar geschlossen, aber jemand machte uns dennoch auf. Nasir unterhielt sich kurz mit dem Mann, der sich an den Fall zu erinnern schien und uns bat, am nächsten Tag wiederzukommen. Wir schauten uns etwas Khoy an und blieben die Nacht dort in einem Apartment. Am Dienstag machten wir uns morgens wieder auf den Weg zum Tazirat und erhielten ein Schreiben. Wir fuhren erneut zum Konditor und holten eine Schachtel Süßigkeiten, diesmal für die Bank. Diese reagierte auf das Geschenk recht schnell, sagte aber, sie bräuchten noch einen Brief vom Hauptgericht. Dort warteten wir wieder. Diesmal vier Stunden! Ich war nie ein geduldiger Mensch, aber in dieser Zeit bin ich es wohl geworden. Um 13:55 Uhr fanden sie im Gericht endlich die richtige Akte. Da Payam und ich dort ohnehin nichts mehr ausrichten konnten, einigten wir uns darauf, dass wir zurück nach Norduz fahren und Nasir die Angelegenheit in Khoy alleine zu Ende führt.

\chapter*{Brilles hämisches Grinsen}
\addcontentsline{toc}{chapter}{Brilles hämisches Grinsen}

Am nächsten Morgen verlassen wir das Hotel in Norduz und machen uns auf den Weg zur Grenze. Unser Ziel ist das Büro des schnurrbärtigen Zollchefs: Wir wollen ihn davon überzeugen, mein Carnet de Passage endlich abzustempeln. Streng genommen hat der Mercedes den Iran bereits verlassen, da er sicher verwahrt hinter der ersten Schranke auf dem Grenzgelände steht. Payam übernimmt wie gewohnt die Vermittlung, während ich dem Klang der persischen Worte lausche, ohne deren Inhalt greifen zu können. Als wir den Raum nach einer gefühlten Ewigkeit verlassen, berichtet mir Payam von einer überraschenden Wende. Der Zollchef hat versichert, dass sie keine Berufung gegen das kommende Gerichtsurteil einlegen werden. Das bedeutet, dass ich nicht die üblichen zwanzig Tage abwarten muss, bis das Urteil rechtskräftig ist und ich das Auto ausführen darf. Sobald das Dokument vorliegt, wollen sie dem Richter ein Schreiben schicken, damit die Angelegenheit noch am selben Tag abgeschlossen werden kann. In diesem Moment durchströmt mich eine Euphorie, die ich seit Wochen nicht mehr gespürt habe. Die reale Aussicht, den Iran tatsächlich mit meinem Mercedes zu verlassen, hatte ich innerlich fast schon aufgegeben. Doch wie es dieses Land so an sich hat, folgt auf jeden Hoffnungsschimmer unmittelbar der nächste Schlag in die Magengrube. Payam erfährt beiläufig, dass die Strafgebühren für Paragraf 109, die Überschreitung der Aufenthaltsdauer für Fahrzeuge, drastisch erhöht wurden. Kostete ein Tag Verspätung früher etwa 1,50 Euro, verlangt das System nun 800.000 Toman, was laut Payam etwa 40 Euro pro Tag entspricht. Auch ohne das alles nachzurechnen, ist mir sofort klar, dass diese Reise eine finanzielle Katastrophe bleibt. Selbst wenn ich vom Schmuggelvorwurf freigesprochen werde, bleibt diese horrende Strafe bestehen. Meine Euphorie ist augenblicklich verflogen. Mitten in diese Ernüchterung platzt der Anruf von Nasir: Er hat es tatsächlich geschafft und die Kaution in Khoy hinterlegt. Der emotionale Tango reißt mich erneut herum, denn diese Nachricht ist der Schlüssel zu meinem Reisepass. Wir warten etwa eine Stunde, bis Brille an der Grenze eintrudelt. Er beachtet uns kaum und tauscht sich erst einmal ausgiebig mit seinen Kollegen aus. Als wir ihn schließlich abfangen und Payam ihm von der hinterlegten Sicherheit berichtet, reagiert er mit einem trockenen: „Ich glaube dir nicht“. In diesem Augenblick muss ich an den Gedanken denken, dass ein Mensch die Welt nicht sieht, wie sie ist, sondern wie er selbst ist. Brille ignoriert jeden weiteren Klärungsversuch und flüchtet sich zum Assistenten des Zollchefs, mit dem wir bisher ein recht positives Verhältnis pflegten. Payam lässt nicht locker, folgt ihm und erklärt dem Assistenten die Situation: Wir müssen wissen, ob ich mit dem Pass ausreisen kann und wie mein aktueller Status im System vermerkt ist. Brille erkennt schließlich, dass sein kindisches Machtgehabe hier keinen Zuspruch findet, und zeigt sich widerwillig gesprächsbereit. Da die offizielle Bestätigung vom Gericht noch auf dem Weg ist, beordert er einen Sicherheitsbeamten, der meinen Pass nehmen und uns zur Passkontrolle begleiten soll. Es dauert eine weitere Stunde, bis dieser erscheint. Während der Wartezeit grinst Brille mich durch die Scheibe seines Büros selbstgefällig an; ich erwidere sein hämisches Gesicht mit nichts als purer Verachtung. An der Passkontrolle folgt die nächste bürokratische Gewissheit: Mein Visum ist abgelaufen. Man erklärt mir, dass ich erst die Ausländerpolizei aufsuchen muss, um eine Ausreiseerlaubnis zu erhalten. Es ist keine Überraschung, aber zumindest werden die letzten Schritte in Richtung Freiheit nun greifbarer. Der Beamte bringt meinen Pass zurück zu Brille. Wir telefonieren derweil mit dem Gericht, um die Übermittlung der Kautionsbestätigung zu beschleunigen. Als Brille den entscheidenden Anruf schließlich erhält, lässt er uns zur Demonstration seiner verbliebenen Macht weitere 45 Minuten warten. Dann händigt er mir den Pass endlich aus und lacht: „Da hast du deinen Pass endlich wieder“. Ich nehme das Dokument schweigend entgegen und spüre erneut nur tiefe Verachtung für diesen Menschen, der mich in diese Lage gebracht hat und nun glaubt, darüber Witze machen zu können. Wir verlassen das Grenzgelände und verbringen den restlichen Tag in Norduz, während die Freiheit zum Greifen nah scheint.

\chapter*{Zwischen Notaren, Lämmern und Hummer}
\addcontentsline{toc}{chapter}{Zwischen Notaren, Lämmern und Hummer}

Der Donnerstagmorgen begrüßte uns gegen 6 Uhr mit einer Sanftheit, die so gar nicht zu meiner inneren Anspannung passen wollte. Im armenischen Grenzgebiet lag ein Hauch von Frühling in der Luft, warm und sonnig. Eigentlich war der Plan simpel: Wir wollten zu einem Freund von Payam in der Nachbarprovinz Ardabil fahren, um ein frisch geschlachtetes Schaf abzuholen. Doch kaum saßen wir im Wagen, fiel uns auf, dass dieser Ort fünf Stunden entfernt ist und damit absolut nicht auf unserem Weg lag. Wir änderten den Kurs also erst einmal zurück nach Täbris, denn mir fehlte noch ein entscheidendes Puzzleteil für meine Abreise: eine Vollmacht für den Mercedes auf Payams Namen. Ich wusste nicht, ob ich jemals wieder ein Visum für den Iran erhalten würde, und wollte das Auto zur Absicherung auf ihn überschreiben lassen. Mit diesem Vorhaben hatten wir uns bereits vor ein paar Tagen in Khoy auseinandergesetzt, ehrlich gesagt mit nur sehr wenig Erfolg. Sobald die Notare merkten, dass ich kein Farsi spreche, winkten sie ab. Fünf Büros hatten wir erfolglos abgeklappert. Die Zeit rannte und in Täbris erhofften wir uns bessere Chancen. Jedoch schließen die meisten Dienstleister am Donnerstag um 12 Uhr mittags, also mussten wir uns beeilen. In Täbris angekommen, klapperten wir einige Notare ab und beim vierten hatten wir tatsächlich endlich Glück. Uns war klar, dass wir einen Übersetzer bräuchten, also wurde Mr. Translation umgehend informiert und er machte sich direkt auf den Weg. Er brauchte allerdings deutlich länger als gedacht, obwohl der Ort sehr nah von seinem Büro aus war. Das Büro des Notars lag in der Nähe einer Moschee und es stellte sich heraus, dass er zuerst zu der falschen Moschee gegangen war. Als er kurz vor 12 Uhr eintraf, begann eine Diskussion darüber, dass Mr. Translation ja nur auf Englisch übersetzte und nicht auf Deutsch. Einen deutschen Übersetzer gab es aber in der gesamten Provinz nicht. Da die staatlichen Server für die Verarbeitung von notariellen Dokumenten um Punkt 12 Uhr offline gehen, konnten wir das leider nicht mehr abschließen. Kommen Sie am Samstag wieder... Etwas genervt darüber, dass es heute nicht geklappt hat, war ich dennoch erfreut, dass wir endlich eine Möglichkeit für die Vollmacht gefunden hatten. Wie fast jeden Tag, wenn wir in Täbris waren, holten wir Payams Vater von der Arbeit ab, brachten ihn nach Hause und machten uns daraufhin auf den langen Weg nach Ardabil. Die Landschaft, die anfangs noch karg und staubig war, verwandelte sich hinter einem Bergpass plötzlich in ein tiefes, sattes Grün. Diese grüne Landschaft verbesserte direkt die Stimmung. Es ist immer wieder spannend und erstaunlich, wie schnell sich Landschaften ändern können. Gegen 18 Uhr erreichten wir das Dorf des Freundes, eines sehr aktiven, kernigen Typs mit nur wenig hinterbliebenen Zähnen. Bereits bevor wir ankamen, hatten wir darüber gesprochen, die Nacht dort zu verbringen, jedoch hatte ich irgendwie ein besseres Gefühl dabei, wieder zurückzufahren. Ich kann gar nicht genau beschreiben, woran das lag. Payam wollte eher bleiben, weil er schon wirklich den ganzen Tag gefahren war, aber ich konnte ihn nach etwas Hin und Her überreden. Vor allem, weil wir seinem Vater seit Wochen versprochen hatten, mit ihm am nächsten Tag in sein Heimatdorf zu fahren. Bei Abgousht, Kuchen und Tee verbrachten wir ca. zwei Stunden dort, nahmen das Fleisch mit und brachen auf. Nach der ersten Stunde übernahm ich das Steuer. Nach kurzer Zeit schlief Payam ein und ich chauffierte uns nach Täbris. Es hat mir richtig Spaß gemacht, selbst zu fahren, weil es schon eine ganze Weile her war. Um diese Uhrzeit waren die Straßen nicht mehr so befahren, was mir ganz gelegen kam. Auch wenn man tagsüber im Iran unterwegs ist und denkt, dass die Leute verrückt fahren, fängt der Wahnsinn auf den Straßen erst an, sobald es dunkel wird. Da es scheinbar keinerlei Richtlinien für Licht am Auto gibt, gibt es sehr viele Leuchtquellen in allen möglichen Farben, die einen blenden können. Wenn man damit noch klarkommt, gibt es aber noch die Fahrer, die von hinten angerast kommen und mit ihrer Lichthupe ein Stroboskopgewitter imitieren. Dabei ist auch vollkommen egal, ob man Platz machen kann oder nicht. Manchmal habe ich das Gefühl, dass die Iraner ihre extrem positive zwischenmenschliche Lebensweise mit einer möglichst aggressiven Fahrweise kompensieren. Aber nach vier Monaten auf Irans Straßen gewöhnt man sich auch dran. Nach einem Tag, an dem wir 15 Stunden im Auto verbracht hatten, kamen wir um 2 Uhr zurück in Täbris an. Wir schliefen etwas aus und machten uns auf den Weg ins Dorf. Wir kamen bei einem alten Haus in einem eher baufälligen Zustand an. Payams Vater und ein Mann haben ein paar Arbeiten auf dem Hof gemacht. Wir kümmerten uns ums Essen. Auch wenn der Ort irgendwie eine kleine Idylle war, ratterte es in meinem Kopf. Der Punkt mit der Strafe für Paragraf 109 saß mir im Nacken. Ich konnte die neue Summe nicht begleichen. Nach etwas Zeit sprach ich Payam nochmal darauf an und ob wir Möglichkeiten hätten, das zu umgehen oder die Summe zu verringern. Payam rief Saleh an und ich verstand, dass sich aus dem Telefonat etwas entwickelte, was positiv schien. Saleh rechnete nochmal nach. 800.000 Toman sind keine 40 Euro, es sind 4 Euro... Hätte ich auch selbst drauf kommen können. Wie gesagt, war ich mittlerweile in einem Zustand, wo man Dinge eher hinnimmt, als sie zu hinterfragen. Mir fiel ein riesiger Stein vom Herzen. Das war der letzte Punkt, der mich wirklich wurmte. Am nächsten Morgen trafen wir uns als Erstes mit Mr. Translation beim Notar. Die Vollmacht wurde aufgeschrieben, Mr. Translation übersetzte den Inhalt, ich unterschrieb. Das lief ausnahmsweise echt reibungslos. Jetzt kam der letzte kritische Punkt: der Weg zur Ausländerpolizei. Dort angekommen, ging der Ping Pong wieder los. Nachdem wir gefühlt mit so gut wie jedem Mitarbeiter gesprochen hatten, brachte man uns zum Chef, der mir die 20 Tage Verlängerung bewilligte. Er sagte nur, dass ich für die 10 Tage, die ich drüber bin, die Strafe zahlen soll. 4.000.000 Toman sind knapp über 20 Euro. Dafür, dass alle sehr ernst wirkten, war es natürlich mein kleinstes Problem. Am Ende kam allerdings wieder die Ernüchterung. Sie behielten meinen Pass zur Verifizierung für einen Tag dort. Sofort gingen bei mir wieder die Alarmglocken an, denn die Warnung über meine Red Flag spukte mir noch immer im Kopf herum. Dass kurz zuvor bereits zwei Männer in der Schlange vor uns direkt zum Geheimdienst geschickt wurden, machte die Sache nicht besser. Das war absolut kein gutes Zeichen. Payam rief den Anwalt an, damit dieser beim Richter anfragte, ob er das Urteil schreiben würde. Die Zeit wurde knapp, bis ich ausreisen konnte, natürlich am liebsten mit dem Auto. Der Anwalt erreichte aber nichts beim Richter. Das Einzige, was uns blieb, war, am nächsten Morgen nochmals Scarface um Hilfe zu fragen. Abends lag ich im Bett und guckte mir Insta-Stories an. Die meisten Nachrichten beziehe ich von worldpolitics\_daily. Dort sah ich ein Bild, wie an diesem Tag die US-Truppen im Nahen Osten Steak, Pie, Krabbenbeine und Hummer serviert bekamen. Mit der Überschrift „Wer weiß, der weiß“ wurde mir sofort klar, was das zu bedeuten hat. Jedem, der sich ein wenig damit auskennt, bleibt bei diesem Anblick das Herz stehen, da es sich oft um die letzte Mahlzeit vor einem großen Einsatz handelt. Diese Panik, die ich bestimmt seit zwei Wochen nicht mehr verspürt habe, stieg wieder in mir hoch. Es war eindeutig an der Zeit, das Land zu verlassen.

\chapter*{Tourist, Tourist}
\addcontentsline{toc}{chapter}{Tourist, Tourist}

Es ist Sonntag. Um 8:30 Uhr stehen wir in Scarfaces Büro. Es besteht die theoretische Chance, dass er den Richter anruft und das Urteil noch heute gefällt wird. Nach einem kurzen Gespräch ist klar: daraus wird nichts. Ein weiterer bürokratischer Leerlauf. Wir fahren erneut zur Ausländerpolizei. Da wir zu früh dran sind, müssen wir 15 Minuten warten. Ich bin deutlich angespannter als sonst. Jede Minute in diesem Gebäude fühlt sich zäh an. Ich setze mich in den Wartebereich, während Payam die Formalitäten klärt. Plötzlich klingelt eines der Telefone. Es steht nicht direkt am Platz des Beamten, sondern etwas abseits. Es ist das gleiche Telefon, über das gestern die Anweisung kam, dass sich die beiden Männer vor uns beim Geheimdienst vorstellen müssen. Mein Herz rutscht mir in die Hose, während der Beamte am Telefon ist und parallel mit Payam spricht. Nach einer gefühlten Ewigkeit dreht Payam sich um. Er hält meinen Pass in der Hand. Ich darf das Land tatsächlich verlassen. Payam musste allerdings seine eigenen Personalien über das Telefon übermitteln. Wir sind uns beide unsicher, was das für ihn bedeutet. Eines ist mir jedoch klar: Wäre er nicht dabei gewesen, hätte ich den Weg zum Geheimdienst antreten müssen. Die letzte Fahrt nach Norduz vergeht wie im Flug. Ich bin aufgeregt, das Land bald verlassen zu können. Eigentlich ist alles für meine Abreise bereit, doch das schriftliche Urteil lässt weiter auf sich warten. Da mein Carnet in zwei Wochen abläuft, fragen wir ein letztes Mal nach, ob sie es nicht doch vorab stempeln können. Wir schlagen vor, dass sowohl der Wagen als auch das Carnet als Pfand auf dem Gelände bleiben, bis die Angelegenheit offiziell geklärt ist. Die Antwort ist ein klares Nein. Immerhin werde ich nach dem Urteil ein Schreiben erhalten, das bestätigt, wann ich an die Grenze kam und dass das Auto seitdem dort stand. In dieser Situation ist das zumindest ein kleiner Strohhalm. Ich gehe ein letztes Mal zum W124. Während ich letzte Sachen ausräume und anderes verstaue, taucht plötzlich der Besitzer eines österreichischen Offroad-Wagens auf. Er parkt direkt neben mir. Es ist surreal: Über einen Monat habe ich an dieser Grenze verbracht und keinen einzigen Europäer gesehen, und genau in diesem Moment steht jemand aus einem Nachbarland vor mir. Er ist mit einem Guide unterwegs und hat ein ähnliches Problem: Er musste seinen Wagen wegen eines Defekts im Land lassen, wenn auch nicht so lange wie ich. Ich erzähle ihm kurz meine Geschichte. Er starrt mich mit großen Augen an. Ich lächle und sage ihm, dass er glücklicherweise nur die Strafe für Paragraf 109 zahlen muss und keine Anzeige wegen Schmuggels am Hals hat. Ich merke, wie unsicher er ist. Es erinnert mich an mich selbst vor vier Wochen: unwissend darüber, was hier wirklich passiert. Ich muss schmunzeln. Meine Sichtweise hat sich komplett gewandelt. Ich bin abgestumpft. Während er und sein Guide zum Hotel aufbrechen, gehe ich in die andere Richtung. Vor der Passkontrolle hat sich eine lange Schlange gebildet. Das ist ungewöhnlich, normalerweise ist dieser Bereich fast leer. Während ich warte, fixiert mich ein Mann vom Geheimdienst von der Seite. Er winkt mich zu sich. Na toll. In meinem Kopf sehe ich schon das stundenlange Verhör vor mir. "What are you doing in Iran?", fragt er in sehr schlechtem Englisch. "Tourist!", antworte ich. Er versucht nachzuhaken, aber sein Vokabular reicht nicht aus. Ich wiederhole nur: "Tourist, Tourist". Er nimmt meinen Pass und begutachtet etwa eine Minute lang die Stempel meiner Reise. Er sieht mir kurz in die Augen und lässt mich dann gehen. Keine weiteren Fragen. Kein Verhör. Es ist soweit. Ich verlasse den Iran am 22.02.26.
